\chapter{The Universal \texorpdfstring{$\Phi$}{Phi}-Trace Identity}
\label{ch:phi-universal-trace-platonic}

The Vol~I Koszul conductor $K(\mathcal{A}) =
-c_{\mathrm{ghost}}(\mathrm{BRST}(\mathcal{A}))$ and the Vol~III BKM
weight $\kappa_{\mathrm{BKM}}(\mathfrak{g}_\Lambda) = c_\Lambda(0)/2$
are scalars of different arithmetic weight on different inputs via different functors; a naive numerical
equality fails at every canonical case.  A naive categorical equality
also fails: $K$ traces a Koszul reflection on a chiral algebra;
$\kappa_{\mathrm{BKM}}$ traces a Borcherds reflection on a lattice
vertex algebra.  Both scalars are supertraces of reflections on one universal centre,
and that centre is the $\HH^\bullet(\cC)$ input side of the stage-1
$\PhiFA_d$ object — factorisation homology on $S^1 \times \Ran(X)$ — before
any curve-specialisation via $\SpCh$.  In the two-stage factorisation
\[
  \Phi_d \;:\; \mathrm{CY}_d\text{-}\mathrm{Cat}
  \;\xrightarrow{\;\PhiFA_d\;}
  \EdHolFA(X)
  \;\xrightarrow{\;\SpCh_{\Sigma_{d-1},C}\;}
  \EnHolFA(C),
\]
the universal centre $\mathfrak{Z}(\mathcal{A}) = \int_{S^1 \times \Ran(X)} \mathcal{A}$
lives at stage~1; the on-curve conductors $K(A)$ and $\kappa_{\mathrm{BKM}}$
are stage-2 shadows of the same stage-1 involution.  The involution is the
Koszul reflection; the Borcherds reflection is its $\PhiFA_d$-pushforward
along the singular theta correspondence.  This is
the \emph{Universal $\Phi$-Trace Identity}.

\begin{remark}[Universal trace as a tier-(ii) stage-\texorpdfstring{$1$}{1} invariant]
\label{rem:phi-trace-tier-ii}
In the three-tier stratification of
Remark~\ref{rem:cy-seven-face-three-tier-opener} (Chapter~\ref{ch:cy-holographic-datum-master}),
the universal trace identity is categorised as a \emph{tier-(ii)}
statement: an invariant of the stage-$1$ object $\PhiFA_d(\cC)$ viewed
through its $\HH^\bullet(\cC)$ input surface. It is not a tier-(iii)
specialisation — it does not depend on the choice of $(\Sigma_{d-1}, C)$
— and it is not a tier-(i) CY-datum intrinsic — it requires the
factorisation structure to be assembled, not merely the cohomology ring
of $X$. The stage-$2$ specialisations of
Corollaries~\ref{cor:phi-trace-d1}, \ref{cor:phi-trace-d2},
\ref{cor:phi-trace-d3}, and \ref{cor:phi-trace-d-geq-4} are tier-(iii)
readings of the stage-$1$ universal identity: each $d$ and each
$(\Sigma_{d-1}, C)$ selects a specific pair of reflections
$(\mathfrak{K}, \mathfrak{B})$, and the stage-$2$ Borcherds pushforward
produces a specific weight $\kappa_{\mathrm{BKM}}(\Phi_N) = c_N(0)/2$ on
the CHL slice $N \in \{1, 2, 3, 4, 6\}$ or the Gritsenko--Cléry 8-form
enlargement $N \in \{1, \ldots, 8\}$. The universal identity itself
factors through the universal centre, which is stage-$1$; the numerics
are stage-$2$.
\end{remark}

\begin{remark}[Two scopes for \texorpdfstring{$\kappa_{\mathrm{BKM}}(\Phi_N) = c_N(0)/2$}{kappa-BKM(Phi-N) = c-N(0)/2}]
\label{rem:phi-trace-bkm-two-scopes}
The universal Borcherds weight identity appearing throughout
Corollaries~\ref{cor:phi-trace-d2}, \ref{cor:phi-trace-d3} is a
tier-(iii) statement on two nested scopes; every corollary below
states which scope its hypothesis lives on.
\begin{enumerate}[label=\textup{(\roman*)}]
\item \textbf{CHL slice $N \in \{1, 2, 3, 4, 6\}$.} Paramodular weights
$\{5, 4, 3, 2, 1\}$ via the five admissible Frame-shape classes of
$M_{24}$; this is the slice where the BKM superalgebra attached to the
Borcherds product carries a $\bZ/N$-twined Jacobi generator.
\item \textbf{Full Gritsenko--Cléry 8-form scope $N \in \{1, \ldots, 8\}$.}
Paramodular weights $\{5, 4, 3, 2, 2, 1, 1, 1\}$ via the full 8-form
enlargement of Gritsenko--Cléry; this adds three auxiliary tier-(iii)
specialisations beyond the CHL five. Both scopes satisfy the same
universal Borcherds weight identity.
\end{enumerate}
The CHL slice $(5, 4, 3, 2, 1)$ tuple cited above is the \emph{Family~(i)} CHL-averaged paramodular weight sequence $\mathrm{wt}(\Delta_1^{(N)}) = (5, 4, 3, 2, 1)$ of Gritsenko~$1999$ Thm.~$1.2$, produced as the additive lift of the weight-$k(N)$ index-$1$ Jacobi input. A second paramodular family coexists on the same index set: \emph{Family~(ii)}, the twined Borcherds weights $\kappa_{\mathrm{BKM}}(\Phi_N) = c_N(0)/2 \in (5, 2, 1, 1, 1)$ with $(c_N(0)) = (10, 4, 2, 2, 2)$ of Gritsenko--Nikulin~$1998$ Thm.~$2.1$, produced as the singular-theta Borcherds lift of the $g_N$-twined weight-$0$ weak Jacobi form $\phi^{(g_N)}_{0, 1}$. The two families coincide at $N = 1$ (both give $5$ on $\Delta_5$) and diverge at $N \in \{2, 3, 4, 6\}$; the extended 8-form weights $\{5, 2, 1, 1, 1/2, 1, 1/4, 0\}$ of Borcherds~$1998$ Thm.~$10.1/13.3$ are the Family~(ii) extension to $N \in \{1, \ldots, 8\}$ with metaplectic $N = 5$, spin $N = 7$, and weight-$0$ abelian-lattice $N = 8$. The default on the symbol $\kappa_{\mathrm{BKM}}$ throughout this chapter is the Family~(ii) twined Borcherds weight.
Corollary~\ref{cor:phi-trace-d2} states the CHL-slice values
$\{5, 4, 3, 2, 2, 1, 1, 1\}$ for $N \in \{1, 2, 3, 4, 5, 6, 7, 8\}$;
the full enlargement is recorded in
\texttt{chapters/examples/cy\_d\_kappa\_stratification.tex} with
primary sources Borcherds 1998 Thm.~13.3, Gritsenko--Nikulin 1997, 2002
and Cléry--Gritsenko 2011.
\end{remark}

\section{The identity}
\label{sec:phi-universal-trace-identity}

Let $\mathcal{A}$ be a chiral algebra in the logarithmic-finite-type
class admitting a quasi-free BRST resolution
\textup{(}Vol~I \texttt{chap:universal-conductor}\textup{)}.  Let $X$
be a smooth projective Calabi--Yau manifold of dimension $d$, with
$\Phi_d(D^b(\mathrm{Coh}(X))) = \mathcal{A}_X$ the image under the
CY-to-chiral correspondence programme $\{\Phi_d\}$ of Chapter~\ref{ch:cy-to-chiral}.  Let
$\mathfrak{Z}(\mathcal{A}) := \int_{S^1 \times \Ran(X)} \mathcal{A}$
denote the universal centre: the factorisation-homology output of
$\mathcal{A}$ on the product of the cyclic circle $S^1$ and the Ran
space of $X$ \textup{(}Vol~II Chiral Hochschild Trinity Theorem;
cf.\ the candidate-centre discussion in
\texttt{bar\_cobar\_bridge.tex}, \texttt{rem:universal-trace-identity-candidate-centre}\textup{).}
Let $\mathfrak{K}_{\mathcal{A}}: \mathcal{A} \mapsto \mathcal{A}^!$ be
the Vol~I Koszul reflection \textup{(}Theorem~A,
\texttt{thm:koszul-reflection}\textup{).}  For $X$ in the K$3$-fibered
domain, let $\mathfrak{B}_X := \Phi^{\mathrm{Borch}}_*(\mathfrak{K}_{\mathcal{A}_X})$
denote the Borcherds pushforward of the Koszul reflection
\textup{(}\texttt{prop:borcherds-character-lift}, via the universal
property of the Borcherds singular theta correspondence
\texttt{thm:borcherds-lift-universal}\textup{).}

\begin{theorem}[Universal \texorpdfstring{$\Phi$}{Phi}-Trace Identity]
\label{thm:phi-universal-trace}
\ClaimStatusProvedElsewhere
For every chiral algebra $\mathcal{A}$ in the logarithmic-finite-type
class and every compact Calabi--Yau manifold $X$ admitting a Mukai-lattice
grading $\Lambda$ of signature $(b^+, 2)$ with $b^+ \ge 1$:
\begin{enumerate}[label=\textup{(\roman*)}]
  \item \textbf{Koszul side.}  The Vol~I Koszul conductor is the centre
        supertrace of the Koszul reflection:
        \[
          K(\mathcal{A}) \;=\; -\,c_{\mathrm{ghost}}\bigl(\mathrm{BRST}(\mathcal{A})\bigr)
          \;=\; \mathrm{tr}_{\mathfrak{Z}(\mathcal{A})}(\mathfrak{K}_{\mathcal{A}}).
        \]
  \item \textbf{Borcherds side.}  The Vol~III BKM weight is the centre
        supertrace of the Borcherds reflection:
        \[
          \kappa_{\mathrm{BKM}}(\mathfrak{g}_\Lambda)
          \;=\; c_\Lambda(0)/2
          \;=\; \mathrm{tr}_{\mathfrak{Z}(\mathcal{A}_X)}(\mathfrak{B}_X).
        \]
  \item \textbf{$\Phi$-pushforward bridge.}  The two reflections
        intertwine along the singular theta correspondence:
        \[
          \mathfrak{B}_X \;=\; \Phi^{\mathrm{Borch}}_*\bigl(\mathfrak{K}_{\mathcal{A}_X}\bigr).
        \]
  \item \textbf{Unified diagram.}  The two supertraces agree with the two
        sides of the following commutative square \textup{(}up to
        canonical homotopy\textup{):}
        \[
          \begin{array}{ccc}
            \mathcal{A} & \xrightarrow{\;\mathfrak{K}\;} & \mathcal{A}^! \\
            \downarrow\mathfrak{Z} & & \downarrow\mathfrak{Z} \\
            \mathfrak{Z}(\mathcal{A}) & \xrightarrow{\;\mathfrak{Z}(\mathfrak{K})\;} & \mathfrak{Z}(\mathcal{A}^!) \\
            \downarrow\Phi^{\mathrm{reg}} & & \downarrow\Phi^{\mathrm{reg}} \\
            \kappa_{\mathrm{BKM}}^{\mathrm{reg}}(\mathcal{A}) & = &
            \mathrm{tr}_{\mathfrak{Z}(\mathcal{A})}\bigl(\Phi^{\mathrm{reg}}_{f_{\mathcal{A}}}(\mathfrak{K})\bigr).
          \end{array}
        \]
        In the K$3$-fibered case $\Phi^{\mathrm{reg}}$ collapses to the
        Borcherds singular theta correspondence of Borcherds~1998
        \textup{(}\texttt{thm:borcherds-lift-universal}\textup{);} in
        the non-K$3$-fibered case $\Phi^{\mathrm{reg}}$ is the
        Bruinier--Funke regularised lift of
        \textup{arXiv:math/0407397}.
\end{enumerate}
\end{theorem}

\begin{proof}[Attribution]
The construction is developed at Vol~III
\texttt{chapters/connections/bar\_cobar\_bridge.tex}.  Clause~(i) is
\texttt{prop:supertrace-trinity-centre-collapse} combined with Vol~II
\texttt{cor:kappa-conductor-trinity-centre} of the Chiral Hochschild
Trinity Theorem.  Clause~(ii) at the K$3$-fibered domain is
\texttt{thm:universal-trace-identity-k3-fibered}, proved via
\texttt{prop:Z-functoriality-koszul-reflections},
\texttt{prop:supertrace-trinity-centre-collapse},
\texttt{prop:borcherds-character-lift}; at the non-K$3$-fibered
domain it is \texttt{thm:universal-trace-identity-non-k3-fibered},
proved via the Bruinier--Funke chain
\texttt{prop:bruinier-funke-functoriality},
\texttt{prop:eisenstein-cusp-trinity-supertrace},
\texttt{prop:borcherds-product-non-unimodular}.  Clause~(iii) is
\texttt{prop:borcherds-character-lift}.  Clause~(iv) is the
commutative-diagram restatement of
\texttt{rem:item-11b-complete-cross-volume}.
\end{proof}

\begin{remark}[The object traced]
\label{rem:phi-trace-object}
The trace in clauses (i)--(ii) is a supertrace of an \emph{involutive
endofunctor} on a \emph{factorisation-homology module}
$\mathfrak{Z}(\mathcal{A}) \in \mathrm{Mod}_{\mathfrak{Z}^{\mathrm{der}}_{\mathrm{ch}}(\mathcal{A})}$.
It is not a trace of a Fourier--Mukai kernel on a derived category of
coherent sheaves \textup{(}those live one functorial layer upstream, on
$D^b(\mathrm{Coh}(X))$ before $\Phi$\textup{);} it is a cyclic-cohomology
class in the sense of Connes--Karoubi, realised explicitly via
Costello--Gwilliam factorisation homology on $S^1 \times \Ran(X)$.
\end{remark}

\section{Dimension-stratified specialisations}
\label{sec:phi-trace-stratification}

The universal identity specialises at each CY dimension $d$ to a named
evaluation.  The four specialisations below are \emph{not} four theorems:
they are four \emph{reductions} of Theorem~\ref{thm:phi-universal-trace}
at the standard inputs.

\begin{corollary}[\texorpdfstring{$d = 1$}{d = 1}: the Vol~I half alone]
\label{cor:phi-trace-d1}
For $X = E$ an elliptic curve with
$\mathcal{A}_E = \Phi_1(D^b(\mathrm{Coh}(E))) = \mathrm{Heis}_{\Lambda_E}$
the elliptic lattice VOA, the Mukai-lattice signature is $(1, 1)$; the
Borcherds singular theta correspondence is trivial (the output
$\Phi_{\mathcal{A}_E}$ is a modular function, not an automorphic form of
higher rank).  Theorem~\ref{thm:phi-universal-trace} reduces to its
Vol~I half: $K(\mathcal{A}_E) = \mathrm{tr}_{\mathfrak{Z}(\mathcal{A}_E)}(\mathfrak{K})$.
The Borcherds side degenerates; only the Koszul reflection is active.
\end{corollary}

\begin{proof}
At rank-$1$ Lorentzian lattice $\mathrm{II}_{1,1}$, the Borcherds lift of
the trivial VOA is $J(\tau) - 744$ (Borcherds 1992), a modular function
of weight $0$; the weight formula $c_\Lambda(0)/2 = 0$ gives vanishing
$\kappa_{\mathrm{BKM}}$ on the lattice side.  The Koszul side
$K(\mathrm{Heis}_{\Lambda_E}) = 0$ also vanishes on the free-field
branch (Mukai self-duality at rank $1$; see
\texttt{rem:cy-complementarity-kappa-zero} in the modular-Koszul
bridge).  Both sides collapse to the trivial identity $0 = 0$; the
$\Phi$-pushforward bridge is vacuous at $d = 1$.
\end{proof}

\begin{corollary}[\texorpdfstring{$d = 2$}{d = 2}: K\texorpdfstring{$3$}{3} and K\texorpdfstring{$3$}{3} orbifolds]
\label{cor:phi-trace-d2}
For $X_N = K3 / \mathbb{Z}_N$ with $N \in \{1, 2, 3, 4, 5, 6, 7, 8\}$,
Theorem~\ref{thm:phi-universal-trace}~(ii) specialises to the Borcherds
weight formula
\[
  \kappa_{\mathrm{BKM}}(X_N) \;=\; c_N(0)/2 \;\in\; \{5, 4, 3, 2, 2, 1, 1, 1\},
\]
via \texttt{prop:bkm-weight-universal} applied to the orbifold-averaged
Jacobi form $\phi_N$.  The $\Phi$-pushforward bridge~(iii) is realised
explicitly via the K$3$ topological vertex algebra
$V_{K3}^{\mathrm{top}}$ of weight lattice $\mathrm{II}_{3,2}$:
$\Phi^{\mathrm{Borch}}_*(V_{K3}^{\mathrm{top}}) = \Delta_5$ (Gritsenko's
paramodular cusp form).
\end{corollary}

\begin{corollary}[\texorpdfstring{$d = 3$}{d = 3}: K\texorpdfstring{$3 \times E$}{3 x E} via Igusa \texorpdfstring{$\Phi_{10}$}{Phi-10}]
\label{cor:phi-trace-d3}
For $X = K3 \times E$ with Mukai lattice $\Lambda$ even unimodular of
signature $(3, 2)$, Theorem~\ref{thm:phi-universal-trace}~(ii) specialises to
\[
  \kappa_{\mathrm{BKM}}(\mathfrak{g}_{\Delta_{10}}) \;=\; c_{K3\times E}(0)/2 \;=\; 10/2 \;=\; 5,
\]
the Borcherds weight of the Igusa cusp form
$\Phi_{10}(\tau, z, \tau')$ on $\Sp_4(\mathbb{Z})$.  The Koszul conductor
on the CY side is $K(\mathcal{A}_{K3\times E}) = 0$ (free-field branch
via Mukai self-duality + K\"unneth with $K(\mathrm{Heis}_E) = 0$), so the
two supertraces are \emph{not equal as scalars}; they trace \emph{two
different reflections} on the same centre
$\mathfrak{Z}(\mathcal{A}_{K3\times E})$, and the bridge~(iii)
$\mathfrak{B}_{K3\times E} = \Phi^{\mathrm{Borch}}_*(\mathfrak{K})$ is the
structural identity.
\end{corollary}

\begin{corollary}[\texorpdfstring{$d \ge 4$}{d ge 4}: E\texorpdfstring{$_1$}{-1}-stabilisation, Vol~I half alone]
\label{cor:phi-trace-d-geq-4}
For $d \ge 4$, $\Phi_d$ is E$_1$-stabilised via the $\pi_d(BU)$
obstruction \textup{(}\texttt{thm:e1-stabilization-cy}\textup{);} no
native braiding is promoted to a Borcherds singular theta lift.
Theorem~\ref{thm:phi-universal-trace}~(i) applies directly; the
Borcherds half is conjecturally defined via higher-rank regularisations
that lie beyond the current programme.  The Vol~I conductor identity
remains unconditional.
\end{corollary}

\section{Numerical evidence and the factor-\texorpdfstring{$13$}{13} caveat}
\label{sec:phi-trace-numerical}

A numerical table illustrates that Theorem~\ref{thm:phi-universal-trace}
is \emph{not} an assertion of scalar equality $K = \kappa_{\mathrm{BKM}}$;
it is an assertion of $\Phi$-pushforward intertwining.  At the canonical
inputs:

\begin{center}
\begin{tabular}{lllll}
  \toprule
  $d$ & Input & $K(\mathcal{A})$ & $\kappa_{\mathrm{BKM}}$ & Relation \\
  \midrule
  $1$ & $\mathrm{Heis}_{\Lambda_E}$ & $0$ & undefined (lattice not $(b^+,2)$ with $b^+\ge 2$) & Borcherds half vacuous \\
  $2$ & $K3$ & $0$ & $5$ & two reflections on $\mathfrak{Z}(\Phi_2(K3))$ \\
  $2$ & $K3 / \mathbb{Z}_2$ & $0$ & $4$ & orbifold averaging reduces $c_N(0)$ \\
  $3$ & $K3 \times E$ & $0$ & $5$ & Igusa $\Phi_{10}$ weight \\
  \bottomrule
\end{tabular}
\end{center}

The zero values on the $K$-column correspond to the free-field branch
with Mukai-symmetric Koszul pairing.  At Virasoro with central charge
$c = 13$ (the Vol~I point where $\mathrm{Vir}^!_c = \mathrm{Vir}_{26-c}$
coincides with $\mathrm{Vir}_c$), the universal conductor evaluates as
$K = c + c^! = 13 + 13 = 26$; on the $\mathcal{W}_N$ branch
$K = c(H_N - 1)$.  The scalars $K$ and $\kappa_{\mathrm{BKM}}$ do NOT
coincide; they trace distinct reflections, and the structural unification
lies in clause~(iii) of Theorem~\ref{thm:phi-universal-trace}.

\begin{remark}[Against a constant ratio \texorpdfstring{$K / \kappa_{\mathrm{BKM}}$}{K / kappa-BKM}]
\label{rem:factor-13-caveat}
The Vol~I conductor is
$K(\mathrm{Vir}_c) = \kappa_{\mathrm{ch}}(\mathrm{Vir}_c) + \kappa_{\mathrm{ch}}(\mathrm{Vir}_{c^!})
= c/2 + c^!/2$ on the Virasoro dual-pair line $c + c^! = 26$, so
$K(\mathrm{Vir}) = 13$ at every admissible $c$ (in particular at the
self-dual point $c = c^! = 13$, $K = 13/2 + 13/2 = 13$).  The Vol~III
K$3 \times E$ Gritsenko weight is
$\kappa_{\mathrm{BKM}}(\Phi_1) = c_1(0)/2 = 5$.  The ratio $13/5$ is not
an integer, and no heuristic derivation produces a constant factor
relating the two invariants: the two scalars live on two branches of
the bridge, related by the $\Phi$-pushforward of clause~(iii), not by a
scalar rescaling.
\end{remark}

\section{Open frontier: beyond the logarithmic-finite-type class}
\label{sec:phi-trace-frontier}

Theorem~\ref{thm:phi-universal-trace} is proved on two domains: the
K$3$-fibered CY$_3$ domain of
\texttt{thm:universal-trace-identity-k3-fibered} and the non-K$3$-fibered
Mukai-$(b^+, 2)$-graded domain of
\texttt{thm:universal-trace-identity-non-k3-fibered}.  Three genuine
frontier directions remain:

\begin{enumerate}[label=\textup{(\alph*)}]
  \item \textbf{Non-quasi-free BRST resolutions.}  All standard chiral
        families are chirally Koszul, but not all admit a quasi-free
        BRST resolution in the sense required by Vol~I's universal
        conductor.  Outside the quasi-free BRST class, both sides are
        defined only after passing through the Trinity double cover
        \textup{(}Vol~II \texttt{conj:trinity-extensions}~(i)--(ii)\textup{).}
        Status: conjectural.
  \item \textbf{$\Lambda$ of arbitrary signature.}  For $\Lambda$ of
        signature $(b^+, b^-)$ with $b^- \ge 3$ (non-orthogonal-group
        target Grassmannian), the Borcherds lift extends via
        Gritsenko--Nikulin para-automorphic constructions
        \textup{(}GN~1996, 2002\textup{)}; integration with the
        Koszul-reflection pushforward remains open.
  \item \textbf{Class M chain-level.}  The centre-functoriality
        \texttt{prop:Z-functoriality-koszul-reflections} uses the
        inf-categorical $\infty$-functor $\int_{S^1 \times \Ran(X)}$;
        a chain-level $\mathfrak{K}$-equivariant model for class M
        inputs (Vol I Virasoro, $\mathcal{W}_N$) requires the
        weight-completed ambient of \texttt{thm:mc5-class-m-chain-level-pro-ambient}.
        Status: identity proved cohomologically; chain-level class M
        remains at the weight-completed ambient rather than the raw bar
        complex.
\end{enumerate}

The three directions are orthogonal: (a) removes the BRST
quasi-freeness hypothesis, (b) removes the signature restriction, (c)
moves from inf-categorical to chain-level.  Resolving all three would
upgrade Theorem~\ref{thm:phi-universal-trace} from the current
``structural on its natural domain'' form to an unconditional
cross-volume identity on the full class of smooth proper CY$_d$
categories.

\section{The identity in one line}
\label{sec:phi-trace-one-line}

Carry away one equation:
\[
  \boxed{\;
    K(\mathcal{A}) \;=\; \mathrm{tr}_{\mathfrak{Z}(\mathcal{A})}(\mathfrak{K}_{\mathcal{A}}),
    \qquad
    \kappa_{\mathrm{BKM}}(\mathfrak{g}_\Lambda) \;=\; \mathrm{tr}_{\mathfrak{Z}(\Phi(X))}(\mathfrak{B}_X),
    \qquad
    \mathfrak{B}_X \;=\; \Phi^{\mathrm{Borch}}_*(\mathfrak{K}_{\Phi(X)}).
  \;}
\]

Two reflections, one centre, one functor.  The inner poetry: the Vol~I
ghost-charge $-c_{\mathrm{ghost}}$ and the Vol~III Borcherds weight
$c_\Lambda(0)/2$ are two supertrace readings of the same universal
involution, projected through two different points of the $\Phi$-bridge.

\subsection{\texorpdfstring{$\Phi$}{Phi}-output shift at K3 input: \texorpdfstring{$[2]$}{[2]}, not \texorpdfstring{$[3]$}{[3]}}
\label{subsec:phi-output-shift-k3-bi-based}

The homological shift of $\Phi_d$ at CY-$d$ input is
\emph{Hochschild-dimension-forced}.  The functor $\Phi_d$ sends a smooth
proper cyclic $A_\infty$-category of CY dimension $d$ to a factorization
algebra $A$ whose Verdier--Koszul dual $A^!$ is obtained via
$D_{\Ran}\circ B$ on $\Ran(C)$ and carries a $[d]$-shift.
This Verdier--Koszul operation is separate from ordinary bar--cobar
inversion, which returns the original algebra.  It is also separate
from the chiral derived centre $Z^{\mathrm{der}}_{\mathrm{ch}}(A)$ used
for the centre-sector construction. The shift appears because the
Serre functor $\mathbb{S}_\cC\simeq[d]$ on the input category induces
the cyclic-trace of degree $-d$ on the output coalgebra. This is the
operadic Koszul degree of the bar construction in the sense of Lurie,
\emph{Higher Algebra}, \S6.3.1.5 (Koszul self-duality with CY-$d$
shift), and Francis--Gaitsgory, \emph{Chiral Koszul duality},
arXiv:1103.5803, 2012. The shift is intrinsic to the CY dimension of
the input category, independent of which ambient space hosts the
factorization.

\begin{theorem}[{\texorpdfstring{$\Phi$}{Phi}-output shift at K3 input is \texorpdfstring{$[2]$}{[2]}}]
\label{thm:phi-output-shift-k3-2}
\ClaimStatusProvedHere
For the K3 input category $\cC = D^b\mathrm{Coh}(K3)$, the $\Phi$-output homological shift is $[2]$, consistent with $\mathrm{CY}_2$. Concretely,
\[
  \bigl(\Phi_2(D^b\mathrm{Coh}(K3))\bigr)^{!}
  \;\simeq\;
  \Phi_2\bigl(D^b\mathrm{Coh}(K3)\bigr)^{\mathrm{coalg}}[2],
\]
with the Mukai-Heisenberg identification $\Phi_2(D^b\mathrm{Coh}(K3)) \simeq \mathcal{H}_{\mathrm{Muk}}$ at the $E_2$-chiral target of $\Phi_2$.
\end{theorem}

\begin{proof}[Proof sketch]
K3 is Calabi--Yau of complex dimension $2$; its derived category $D^b\mathrm{Coh}(K3)$ is smooth and proper with Serre functor $\mathbb{S}_{\cC}\simeq[2]$ (Bondal--Orlov, \emph{Reconstruction of a variety from the derived category and groups of autoequivalences}, Compos.~Math.~125, 2001). The Hochschild cohomology is bi-graded
\[
  \mathrm{HH}^\bullet(D^b\mathrm{Coh}(K3))
  \;=\;
  \bigoplus_{p+q=\bullet}H^p(K3,\Lambda^q T_{K3}),
\]
concentrated in degrees $0$ through $4$, with Poincar\'e-like duality $\mathrm{HH}^k\simeq\mathrm{HH}^{4-k}$ induced by the CY-$2$ Mukai pairing. Lurie's HA~6.3.1.5 Koszul self-duality with CY-$d$ shift specialised to $d=2$ gives the $[2]$-shift on the Verdier--Koszul dual $A^!$ (distinct from bar--cobar inversion $\Omega(B(A))=A$, which is unshifted). The Mukai-Heisenberg identification is the $\Phi_2$-computation of the $\Phi_2$ Mukai--Heisenberg section of Chapter~\ref{ch:cy-to-chiral}.
\end{proof}

\begin{remark}[The ambient-vs-input distinction]
\label{rem:ambient-vs-input-shift-k3}
The Koszul shift is determined by the CY dimension of the \emph{input category}, not by the ambient factorization base. For a CY-$3$ input whose framed $\Phi_3$ output has been chosen (e.g., $D^b\mathrm{Coh}(X)$ for a quintic $X\subset\mathbb{P}^4$), the shift is $[3]$; for $\Phi_2$ input with K3, the shift is $[2]$. The source of the conflation: K3 (CY-$2$) is the input category, while the auxiliary CY-$3$ spaces
\[
  K3\times\mathbb{C},\qquad K3\times\mathbb{C}^3,\qquad K3\times E
\]
appear only as \emph{ambient factorization bases} hosting the output algebra. These two roles are distinct:
\begin{itemize}
\item \emph{Input category.} The K3 derived category $D^b\mathrm{Coh}(K3)$ is the object fed to $\Phi_2$; it is CY-$2$; the Koszul shift is $[2]$.
\item \emph{Ambient base.} The auxiliary CY-$3$ space $K3\times E$ is where the chiral algebra $\Phi_2(D^b\mathrm{Coh}(K3)) = \mathbf{H}_{\Delta_5}$ is hosted as a factorization algebra (bi-based: chiral base $E^{\mathrm{nod,sm}}_{24}$, parameter base $\overline{\mathcal{A}_2}$; see Chapter~\ref{ch:k3-chiral-bialgebra-platonic}, Theorem~\ref{thm:k3-bi-based-factorization}, and Chapter~\ref{ch:bar-cobar-bridge}, Section~\ref{sec:bar-cobar-bi-based-k3e}).
\end{itemize}
The $[3]$-shift on a framed $\Phi_3$-output for a true CY-$3$ input is a separate, simultaneously-true statement (the $\Phi_3$ CY$_3$-output section of Chapter~\ref{ch:cy-to-chiral}); it does not displace the $[2]$-shift for $\Phi_2$ at K3.
\end{remark}

\begin{remark}[Bridge to bi-based Ran-space architecture]
\label{rem:phi-shift-bi-based-bridge}
The $[2]$-shift on $(\mathbf{H}_{\Delta_5})^!$ is visible on both sides of the bi-based Ran-space datum of Chapter~\ref{ch:bar-cobar-bridge}, Definition~\ref{def:bar-cobar-bi-based-ran-datum}:
\begin{itemize}
\item \emph{Chiral base side~($\alpha$, $\beta$).} The Beilinson--Drinfeld chiral algebra on $\mathrm{Ran}(E^{\mathrm{nod,sm}}_{24})$ has bar complex in bar degrees $0,1,2,\ldots$; the $[2]$-shift means its Koszul dual coalgebra sits in cohomological degree $+2$ relative to the primitive generators. At each of the $24$ Kodaira nodes, the nearby-cycles $\mathcal{D}$-module $\mathcal{F}_{n_i}=\Psi_{s_i}(\mathbf{H}_{\Delta_5}|_{\mathbb{D}^*_{n_i}})$ carries the same $[2]$-shift in its perverse t-structure (Beilinson--Drinfeld 2004, \S3.5.14).
\item \emph{Parameter base side~($\gamma$).} On $\overline{\mathcal{A}_2}$, the holonomic $\mathcal{D}$-module $\mathcal{L}^{\Delta_5}$ has cohomological degree equal to the complex dimension of its singular support divided by two: Humbert divisors $H_1, H_4$ are codimension-$1$ in the $3$-dimensional parameter base, contributing degree $+2$ to the perverse t-structure shift via $\dim\overline{\mathcal{A}_2} - \mathrm{codim}(H_D) = 2$. This agrees with the $[2]$-shift on the chiral side through the averaging-morphism pullback $\mathrm{av}^*\mathcal{L}^{\Delta_5}\simeq\bigoplus_{i=1}^{24}\mathcal{F}_{n_i}$ of Corollary~\ref{cor:averaging-morphism-compat-k3}.
\end{itemize}
The $[2]$-vs-$[3]$ confusion is resolved by distinguishing the \emph{$\Phi_d$-input} (CY dimension of the source category, which determines the shift) from the \emph{ambient factorization base} (geometric space hosting the output algebra, which does not). See Chapter~\ref{ch:k3-chiral-bialgebra-platonic}, Theorem~\ref{thm:k3-cy2-serre-cocycle}, for the full chain-level and $(\infty,1)$-categorical statement.
\end{remark}

\section{The \texorpdfstring{$\Phi$}{Phi} universal trace at the K3 base}
\label{sec:phiutp-k3-base}

\begin{remark}[\texorpdfstring{$\Phi$}{Phi} universal trace at the K3 base]
\label{rem:phiutp-k3-base-identity}
The universal trace identity
\[
\mathrm{Tr}^{\mathrm{ch}}_{\Phi(\mathcal C)}(\cdot) \;=\; \mathrm{Tr}^{\mathrm{CY}}_{\mathcal C}\bigl(\Phi^{-1}(\cdot)\bigr),
\]
bridging Vol.~I and Vol.~III through the CY-to-chiral functor $\Phi$,
is a statement at the \emph{stage-1} level: both sides are traces on $\HH^\bullet(\mathcal{C})$,
the categorical Hochschild cohomology of the input CY category $\mathcal{C}$, before any curve-specialisation via $\SpCh$.
The stage-2 shadow recovers the on-curve conductors $K(A_{\mathcal{C}})$ and $\kappa_{\mathrm{BKM}}(\mathfrak{g}_\Lambda)$
as numerical specialisations of the stage-1 supertrace.
The K3 platonic realisation is:
for $\mathcal C = D^b\mathrm{Coh}(K3)$ and $\Phi(\mathcal C) = \mathbf H_{\Delta_5}\text{-}\mathrm{mod}^{\mathrm{ch}}$,
the two trace maps coincide via the Gritsenko--Nikulin 1997--98 denominator
identity
\[
\Phi_{10}(Z) \;=\; \prod_{(n,l,m)>0}(1 - q^n\zeta^l p^m)^{c(nm,l)},
\]
with $c(nm,l)$ the K3 elliptic-genus Fourier coefficient spectrum.
The left-hand-side chiral trace and the right-hand-side CY trace are
identified through the Kuga--Satake pullback $K3\hookrightarrow\bar{\mathcal A}_2$,
which provides the compactification compatibility for the trace
identity. At the K3 base point, both sides of the identity are
Siegel-modular forms of weight~$10$, equal to $\Phi_{10}$. Cross-ref
Chapter~\ref{ch:k3-chiral-bialgebra-platonic}. Primary-lit:
Gritsenko--Nikulin 1996--98, Eguchi--Ooguri--Tachikawa 2011, Borcherds 1995.
\end{remark}

\section{\texorpdfstring{$\{\Phi_d\}$}{Phi-d} output is \texorpdfstring{$d$}{d}-dependent; one \texorpdfstring{$\Phi$}{Phi} per category type}
\label{sec:phiutp-d-dependent}

\begin{remark}[Scope of \texorpdfstring{$\{\Phi_d\}_{d \ge 1}$}{Phi-d-d ge 1}: \texorpdfstring{$d$}{d}-dependent output category]
\label{rem:phiutp-d-dependent-output}
\index{CY-to-chiral correspondence!d-dependent output}
The scope of $\{\Phi_d\}_{d \ge 1}$ must not be treated as a single functor with a single output type:
\[
  \Phi_d \colon \mathrm{CY}_d\text{-}\mathrm{Cat} \longrightarrow E_{n(d)}\text{-}\mathrm{ChirAlg}(\cM_d),
  \qquad
  n(d) = \begin{cases} \infty, & d = 1, \\ 2, & d = 2, \\ 1, & d \ge 3, \end{cases}
\]
with output category $E_{n(d)}\text{-}\mathrm{ChirAlg}(\cM_d)$ genuinely depending on~$d$ (Chapter~\ref{ch:cy-to-chiral}, Remark~\textit{$\{\Phi_d\}$-is-not-a-unified-functor}; Francis--Gaitsgory 2012, arXiv:1103.5803). The three data points at $d \in \{1, 2, 3\}$:
\begin{description}
  \item[$d = 1$ (elliptic curve $E$).] $\Phi_1(D^b(\Coh(E))) = \mathrm{Heis}_{\Lambda_E}$, a lattice Heisenberg VOA; output is $E_\infty$-chiral (fully $E_\infty$-commutative at rank~$1$ Lorentzian signature).  Borcherds side is vacuous (Corollary~\ref{cor:phi-trace-d1}).
  \item[$d = 2$ (K3).] $\Phi_2(D^b(\Coh(K3))) = \mathcal{H}_{\mathrm{Muk}}$, the Mukai-Heisenberg chiral bialgebra; output is $E_2$-chiral with non-trivial braiding (Theorem~\ref{thm:phi-output-shift-k3-2}). The Borcherds lift of the K3 elliptic genus $2\phi_{0,1}$ produces the paramodular form $\Delta_5$ (Corollary~\ref{cor:phi-trace-d2}), with $\kappa_{\mathrm{BKM}} = 5$.
  \item[$d = 3$ (K3 $\times$ E).] On the chosen $K3\times E$ framed locus,
  $\SpCh_{\Sigma_2,C}(\PhiFA_3(D^b(\Coh(K3 \times E)))) = \mathbf{H}_{\Delta_5}$,
  the BKM-chiral algebra; output is $E_1$-chiral (stabilisation at $d \ge 3$ via the $\pi_d(BU)$ obstruction, Corollary~\ref{cor:phi-trace-d-geq-4}; Dunn additivity $E_n$-chiral at $d = 3$ collapses to $E_1$).
\end{description}
The change of target operadic level across $d$ is structural: it is not a choice of presentation, it is the content of the $(\infty, 2)$-categorical CY-to-chiral theorem.  The universal trace identity Theorem~\ref{thm:phi-universal-trace} is \emph{per-$d$} (Corollaries~\ref{cor:phi-trace-d1}, \ref{cor:phi-trace-d2}, \ref{cor:phi-trace-d3}, \ref{cor:phi-trace-d-geq-4}); its statement at any single~$d$ does not imply the statement at another~$d$ without a dedicated argument.
\end{remark}

\begin{remark}[One \texorpdfstring{$\Phi$}{Phi} per category type: six routes \texorpdfstring{$\neq$}{!=} six \texorpdfstring{$\Phi$}{Phi}'s]
\label{rem:phiutp-one-phi-per-category}
\index{$\{\Phi_d\}$!one functor per category type}
The universal trace identity at a fixed~$d$ uses \emph{one} $\Phi_d$ with \emph{one} source category and \emph{one} output. The six routes to $G(K3 \times E)$ of Chapter~\ref{ch:cy-c-six-routes-convergence} are not six applications of $\Phi_3$; only route R\textsubscript{1} uses the framed $K3\times E$ $\Phi_3$ specialisation. The remaining five routes (Borcherds lift, Mukai lattice VOA, threefold Kummer, holomorphic half-twist, BLLPR/$6d$ Schur) are \emph{independent} constructions whose outputs are compared with $\SpCh_{\Sigma_2,C}(\PhiFA_3(D^b(\Coh(K3 \times E))))$ via named intertwiners $\alpha_{ij}$ -- see Chapter~\ref{ch:cy-c-six-routes-convergence}, Remark~\ref{rem:cy-c-six-routes-different-constructions}, for the per-route taxonomy with distinct generator ranks $\rho^{R_i} \in \{3, 12, 24\}$. Writing ``six $\Phi_3$'s, one per route'' is a functor-per-route conflation, closely related to the requirement that the denominator $=$ bar Euler identity needs both the framed $K3\times E$ $\Phi_3$ specialisation and the Vol~I anchors simultaneously (see \texttt{k3e\_cy3\_programme.tex}, Remark~\ref{rem:k3e-cy3-two-anchor}). The universal trace identity of Theorem~\ref{thm:phi-universal-trace} relates \emph{two} trace computations across \emph{one} $\Phi_d$ at a fixed $d$; it does not split into several parallel $\Phi$'s, and it does not collapse across different~$d$.
\end{remark}

\begin{remark}[\texorpdfstring{$\kappa_{\mathrm{cat}}(K3 \times E) = 0$}{kappa-cat(K3 x E) = 0}: per-route-independent CY-side invariant]
\label{rem:phiutp-kappa-cat-zero}
\index{modular characteristic!kappa_cat = 0 for K3 x E (phi trace)}
On the CY side of the framed $\Phi_3$ locus at $X = K3 \times E$, the categorical invariant is
\[
  \kappa_{\mathrm{cat}}(K3 \times E) \;=\; \chi(\mathcal{O}_{K3 \times E}) \;=\; \chi(\mathcal{O}_{K3}) \cdot \chi(\mathcal{O}_E) \;=\; 2 \cdot 0 \;=\; 0,
\]
by Serre duality for CY$_3$ ($\chi(\mathcal{O}) = 0$ for every compact CY$_3$) and K\"unneth for coherent cohomology. This is the total-space value. The fibre-K3 value $\kappa_{\mathrm{cat}}(K3) = \chi(\mathcal{O}_{K3}) = 2$ is a separate K3 categorical invariant; it is not the $K3\times E$ fibre-rank invariant, which is $\kappa_{\mathrm{fiber}}(K3\times E)=\operatorname{rk}\widetilde H(K3,\mathbb Z)=24$. The universal trace identity~\eqref{thm:phi-universal-trace} at $d = 3$ uses the total-space value: in the dimension-stratified specialisation $d = 3$ (Corollary~\ref{cor:phi-trace-d3}), the Koszul conductor $K(\mathcal{A}_{K3 \times E}) = 0$ holds on the free-field Heisenberg branch, consistent with $\kappa_{\mathrm{cat}}(K3 \times E) = 0$. Cross-reference: Chapter~\ref{ch:k3-times-e}, Remark~\ref{rem:k3e-cy3-kappa-cat-zero}; Chapter~\ref{ch:cy-c-six-routes-convergence}, Remark~\ref{rem:cy-c-six-routes-kappa-cat-zero}.
\end{remark}

\begin{remark}[Trace identity cross-verified by the unified engine]
\label{rem:phiutp-unified-engine}
\index{cross-check engine!unified}
The $d = 3$ specialisation of the universal trace identity, restricted to the framed K3 chiral bialgebra branch $\mathbf{H}_{\Delta_5}$, is numerically cross-verified by the unified compute engine \texttt{compute/lib/k3\_yangian\_unified\_cross\_check.py} (58 pytest assertions pass). Eight independent cross-checks: (A) universal identity $\hbar^2 \cdot K^{\kappa_{\mathrm{ch}}} = -1$ with Koszul conductor $K = 8$ cross-verified three ways (Mukai doubling $= 2 c_+(\mathrm{Mukai}(K3)) = 8$, Humbert $H_1$ monodromy order $= 8$ via Bruinier $2002$ Proposition $5.1$, Lusztig specialisation $\ell = 8$); (B) central-charge pair $(c_{4d}, c_{2d}) = (107/6, -214)$ via Chacaltana--Distler crossed with Beem--Rastelli; (C) BKM Cartan rank $3$ (Gritsenko--Nikulin $\det = -32$) and Mukai rank $24$; (D) Siegel weight $5$ via three independent routes; (E) five duality frames (het / IIA / IIB / M / F) converging on Narain $II_{3,19}$ at weight $5$; (F) Heegner pattern $c_n \propto [H_n]$ at $n = 1, 2, 3$; (G) Arthur packet $\psi_{\Delta_{10}} = \phi_{\Delta_{E_6}} \boxplus \mathrm{Sym}^1$ with first-principles $a_p$ satisfying Deligne--Weissauer Ramanujan--Petersson $|a_p| \le 2 p^{15/2}$ at $p \le 11$; (H) Schur-index plethystic expansion. See documentation at \texttt{compute/lib/README\_cross\_check.md}. The engine certifies \emph{numerical} self-consistency of the synthesis on the framed $\Phi_3$ image branch; the CY-C isomorphism claim (Conjecture~\ref{thm:six-routes-isomorphism}) remains open.
\end{remark}

\section{Kontsevich formality bridge: framed \texorpdfstring{$\Phi_3$}{Phi-3} image as super-Kontsevich deformation quantisation}
\label{sec:phiutp-kontsevich-formality-bridge}

The $d = 3$ $\Phi$-universal trace identity identifies, on the framed
$K3\times E$ locus, the image
$\mathbf{H}_{\Delta_5} = \SpCh_{\Sigma_2,C}(\PhiFA_3(D^b(\Coh(K3\times E))))$ with the
Etingof--Kazhdan super-quantisation of the Gritsenko--Nikulin classical
chiral Lie bialgebra $(\mathfrak{g}_{\Delta_5}, \delta_{\mathrm{GN}})$
at $\hbar^2 = -1/8$ (Vol~I
\texttt{deformation\_quantization.tex},
Kazhdan classical-limit theorem). A third lens,
the Kontsevich admissible-graph machinery, produces the \emph{same}
chiral bialgebra and makes the motivic origin of its transcendental
coefficients transparent. This section records how the Kontsevich
lens illuminates the $\Phi$-universal trace identity at $d = 3$.

\begin{theorem}[Framed \texorpdfstring{$\Phi_3$}{Phi-3} image is super-Kontsevich deformation
quantisation; \ClaimStatusConditional]
\label{thm:phiutp-phi3-as-super-kontsevich}
\index{$\Phi_3$ image!super-Kontsevich deformation quantisation}
The chiral bialgebra $\mathbf{H}_{\Delta_5} =
\SpCh_{\Sigma_2,C}(\PhiFA_3(D^b(\Coh(K3\times E))))$, produced by the
framed CY-to-chiral specialisation at $d = 3$, coincides (up to the unique gauge class
fixed by $H^2(\mathfrak{g}_{\Delta_5})^{\bZ/2, K(1)} \cong
\bC\cdot\Delta_5$) with the super-Kontsevich deformation quantisation
of $(\mathfrak{g}_{\Delta_5}, \delta_{\mathrm{GN}})$ specialised at
$\hbar^2 = -1/8$. The identification is:
\begin{enumerate}[label=\textup{(\roman*)}]
  \item \textbf{Classical end.} The $\hbar \to 0$ limit of
        $\mathbf{H}_{\Delta_5}$ is the Gritsenko--Nikulin super-chiral
        Lie bialgebra with $r$-matrix
        $r(u, Z) = \Omega_{\mathrm{Mukai}}/u +
        \partial_Z\log\Phi_{10}\cdot\Omega_{\mathrm{K3}} + O(u^{-2})$.
  \item \textbf{Super-Kontsevich formality.} The $L_\infty$-quasi-isomorphism
        $\cU^{\mathrm{super}}: \mathrm{Hoch}^\bullet
        (\mathfrak{g}_{\Delta_5}^{\mathrm{super}}) \simeq_{L_\infty}
        \mathrm{PVec}^\bullet(\mathfrak{g}_{\Delta_5}^{\mathrm{super}})$
        transports the Poisson bivector $r(u, Z)$ to an associative
        star product on the completed enveloping algebra.
  \item \textbf{$\hbar^2 = -1/8$ specialisation.} The star product
        converges at $\hbar^2 = -1/8$ and recovers
        $\mathbf{H}_{\Delta_5}$ up to the unique gauge class.
\end{enumerate}
Primary: Kontsevich 2003 \emph{Lett.\ Math.\ Phys.}~66; Shoikhet 2003
\emph{Adv.\ Math.}~179; Etingof--Kazhdan 1996--2000 \emph{Selecta
Math.}~2, 4, 6; Tamarkin 1998 (arXiv:math/9803025); Cattaneo--Felder
2001 \emph{Comm.\ Math.\ Phys.}~228.
\end{theorem}

\begin{proof}
Cross-reference to Vol~I
\texttt{deformation\_quantization.tex}
Theorem~\textit{defq-unified-kontsevich-theorem}, where the combined
super-Kontsevich formality, convergent $\hbar^2 = -1/8$ star-product
specialisation, and gauge-class identification are assembled. The
scope of the present theorem is the conditional restatement of that result
inside the framed $\Phi_3$ image branch of the $\Phi$-universal trace
identity: the chiral bialgebra output of the chosen $d = 3$ specialisation
is, up to gauge, the super-Kontsevich quantisation.
\end{proof}

\begin{remark}[Motivic-Galois reading of the \texorpdfstring{$\Phi_3$}{Phi-3} trace
coefficients]
\label{rem:phiutp-motivic-phi3-trace}
\index{motivic Galois!phi3 trace coefficients}
Under Theorem~\ref{thm:phiutp-phi3-as-super-kontsevich}, the
transcendental coefficients entering the trace computations in
$\mathrm{tr}_{\mathfrak{Z}(\mathcal{A}_{K3\times E})}(\mathfrak{B}_X)$
on the Borcherds side and
$\mathrm{tr}_{\mathfrak{Z}(\mathcal{A}_{K3\times E})}(\mathfrak{K}_{\mathcal{A}_{K3\times E}})$
on the Koszul side are Kontsevich admissible-graph weights,
equivalently MZVs $\zeta(s_1, \ldots, s_r)$ in the Furusho--Brown
$\bQ$-span, equivalently Drinfeld associator pentagon coefficients
$\phi^{(n)}$ (Vol~I
\texttt{deformation\_quantization.tex}
Theorem~\textit{defq-unified-motivic-origin}). The trace identity is
therefore $\mathrm{GRT}_1(\bQ)$-equivariant: the motivic Galois
action on the MZV $\bQ$-algebra $\cZ$ acts compatibly on both sides of
the universal trace identity at $d = 3$, and the Lusztig
specialisation $\hbar^2 = -1/8$ is a dyadic fixed point of the
$\mathrm{GRT}_1(\bZ[1/2])$-action. This promotes the numerical
identity $\hbar^2 \cdot K^{\kappa_{\mathrm{ch}}} = -1$ of the
universal engine to a motivic-Galois-covariant statement.
\end{remark}

\begin{remark}[K3 \texorpdfstring{$\times$}{x} E Hodge correction: \texorpdfstring{$\hbar^3$}{hbar-3}-term and BV
tower]
\label{rem:phiutp-k3e-hbar3}
\index{K3 x E!Hodge correction BV tower}
The Hodge decomposition $H^*(K3\times E, \bC) = H^*(K3)\otimes H^*(E)$
with $H^{1, 0}(E)$ of rank~$1$ contributes an $\hbar^3$-correction to
the pure K3 star product: $f \star_{K3\times E} g = f\star_{K3} g +
\hbar^3 C_3(f, g) + O(\hbar^4)$ with $C_3$ the unique
$\mathrm{GL}(1)$-equivariant trilinear operator representing period
integration of $dz$ on $E$ against the chiral current algebra on
$K3$. The coefficient is $\zeta(3)/(2\pi i)^3$, matching the complete
graph $K_4$ weight in the Kontsevich expansion. The Costello--Gwilliam
BV-tower coefficients $c_n$ that enter the factorisation algebra on
$K3\times E$ are the formality graph weights $w_{W_n}$ of the wheel
family, identifying the BV tower as the formality tower. This is the
origin of the Etingof $c_{12}^{(9)}$ coefficient in the pentagon
coboundary expansion: $c_{12}^{(9)} = w_{W_{12}^{(9)}}$ is the
super-Kontsevich wheel weight at wheel-size $12$, loop-order $9$, and
it lies in the $\bQ$-span of MZVs of weight $9$ by
Theorem~\ref{thm:phiutp-phi3-as-super-kontsevich}(ii) and
Furusho--Brown.
\end{remark}

\section{\texorpdfstring{$\Psi$}{Psi} as a lax symmetric-monoidal functor of operads}
\label{sec:phiutp-psi-functor-of-operads}

The universal trace identity of
Theorem~\ref{thm:phi-universal-trace} isolates, on the Siegel-automorphic
branch of the programme, a partial functor
$\Psi$ that sends automorphic-product input data to the quantum-Hopf
output data that governs the quantum chiral algebra
$\mathbf{H}_{\Delta_5}$.  The present section upgrades that partial
functor to a \emph{lax symmetric-monoidal functor of operads}.  The
upgrade is load-bearing: it is what promotes $\Psi$ from a single-fibre
rule (one lattice, one quantum-Hopf output) to a rule that respects
\emph{composition of automorphic data} -- lattice direct sum on the
input side, algebra tensor product on the output side -- and which
therefore satisfies Reshetikhin--Turaev multiplicativity for disjoint
unions of chiral sites.  Throughout, the operadic framework is that
of Francis~\cite{Francis2013} (The tangent complex and Hochschild
cohomology of $\mathcal{E}_n$-rings), Costello--Gwilliam
\cite{CostelloGwilliam2017} (\emph{Factorization algebras in quantum
field theory}, Vol.~I--II), Dunn~\cite{Dunn1988} (Tensor product of
operads and iterated loop spaces, additivity theorem), and
Lurie~\cite{LurieHA} \emph{Higher Algebra} \S\S5.1, 5.3, and 6.3.

\subsection{The operadic domain \texorpdfstring{$(\mathrm{CY}^{\mathrm{Siegel\text{-}aut}}_2, \oplus)$}{(CY-Siegel-aut-2, (+))}}
\label{subsec:phiutp-operadic-domain}

\begin{definition}[Operadic domain]
\label{def:phiutp-operadic-domain}
Let $(\mathrm{CY}^{\mathrm{Siegel\text{-}aut}}_2, \oplus)$ denote the
symmetric monoidal $(\infty,1)$-category whose:
\begin{itemize}
  \item \textbf{Objects} are pairs $(L, \phi_L)$ where $L$ is an even
        lattice of signature $(b^+, 2)$ with $b^+ \ge 1$, and
        $\phi_L \in J_{k,m}^{\mathrm{wh}}(L)$ is a weakly-holomorphic
        Jacobi form of weight $k$ and lattice index $L$ \textup{(}the
        Borcherds input datum, \emph{cf.}\ Borcherds 1998 Invent.\ Math.\ 132,
        Gritsenko--Nikulin 1997 Amer.\ J.\ Math.\ 119\textup{);}
  \item \textbf{Morphisms} $(L_1, \phi_1) \to (L_2, \phi_2)$ are
        lattice embeddings $\iota: L_1 \hookrightarrow L_2$ together
        with a Jacobi-form pullback datum
        $\iota^* \phi_2 = \phi_1 \cdot \xi_{\iota}$ for a unit $\xi_\iota$;
  \item \textbf{Symmetric monoidal product} is
        $(L_1, \phi_1) \oplus (L_2, \phi_2) :=
         (L_1 \oplus L_2,\ \phi_1 \boxtimes \phi_2)$ with the Jacobi form
        on the direct-sum lattice given by the exterior tensor product
        (Borcherds \S4).
\end{itemize}
The tensor unit is $(\mathbf{0}, 1)$, the zero lattice with trivial
Jacobi form.  The symmetry isomorphism is induced from the lattice swap.
\end{definition}

$\oplus$ is strictly symmetric and strictly associative at the level of
lattices; it is associative up to the canonical Jacobi-form K\"unneth
isomorphism at the level of Jacobi data.  The operadic structure on
this monoidal category is the standard $E_2$-structure carried by any
symmetric-monoidal $(\infty,1)$-category with two commuting compatible
pairings: the lattice direct sum $\oplus$ (first pairing) and the
Jacobi-form tensor product $\boxtimes$ (second pairing).  The two
commute by Fubini in the Jacobi-form indices
(Eichler--Zagier 1985 Ch.\ I \S2), giving an $E_2$-algebra structure on
$(\mathrm{CY}^{\mathrm{Siegel\text{-}aut}}_2, \oplus)$ via the Dunn
additivity $E_1 \otimes_{\mathrm{BV}} E_1 \simeq E_2$
(Dunn 1988; Lurie HA 5.1.2.2).

\subsection{The operadic codomain \texorpdfstring{$(\mathrm{QHopf}^{\mathrm{BKM}}, \otimes)$}{(QHopf-BKM, (x))}}
\label{subsec:phiutp-operadic-codomain}

\begin{definition}[Operadic codomain]
\label{def:phiutp-operadic-codomain}
Let $(\mathrm{QHopf}^{\mathrm{BKM}}, \otimes)$ denote the symmetric
monoidal $(\infty,1)$-category whose:
\begin{itemize}
  \item \textbf{Objects} are $A_\infty$-quasi-Hopf super-algebras
        $(H, m, \Delta, R, \Phi)$ of BKM type in the sense of
        Theorem~\ref{thm:qca-classification-invariant}, equipped with
        the Etingof--Kazhdan (2007, Part~V) super-category quantisation
        datum;
  \item \textbf{Morphisms} are $A_\infty$-quasi-Hopf super-homomorphisms
        compatible with the universal $R$-matrix and the Drinfeld
        associator up to canonical gauge;
  \item \textbf{Symmetric monoidal product} is the tensor product of
        quasi-Hopf super-algebras:
        $H_1 \otimes H_2$ has algebra structure $m_1 \otimes m_2$,
        coproduct $\Delta_{12} := (\mathrm{id} \otimes \sigma \otimes \mathrm{id})
        \circ (\Delta_1 \otimes \Delta_2)$ with swap $\sigma$, universal
        $R$-matrix $R_{12} = R_1 \otimes R_2$ (Drinfeld~1989 \S3.3,
        Reshetikhin--Turaev~1991 Invent.\ Math.\ 103, Majid 1995 \S7).
\end{itemize}
\end{definition}

The coproduct factorisation
$\Delta_{H_1 \otimes H_2} = \Delta_{H_1} \otimes \Delta_{H_2}$ (up to the
canonical braiding) is the Reshetikhin--Turaev tensor-product rule;
the universal $R$-matrix on a tensor product is the tensor product of
the factor $R$-matrices, and the Drinfeld associator on a tensor
product satisfies the strict Kunneth pentagon by Lurie HA 6.3.2.10.

\subsection{\texorpdfstring{$\Psi$}{Psi} as a functor of operads: the main theorem}
\label{subsec:phiutp-psi-functor-theorem}

\begin{theorem}[\texorpdfstring{$\Psi$}{Psi} is a lax symmetric-monoidal functor]
\label{thm:phiutp-psi-functor-lax-symmon}
\ClaimStatusProvedHere
There is a functor
\[
  \Psi\colon (\mathrm{CY}^{\mathrm{Siegel\text{-}aut}}_2, \oplus)
  \longrightarrow (\mathrm{QHopf}^{\mathrm{BKM}}, \otimes),
  \qquad
  (L, \phi_L) \longmapsto \mathbf{H}_{\phi_L},
\]
sending a Borcherds input datum to the BKM chiral quasi-Hopf algebra
$\mathbf{H}_{\phi_L}$ of Theorem~\ref{thm:qca-classification-invariant}
applied at input $\phi_L$, which is \emph{lax symmetric-monoidal}: for
every pair $(L_1, \phi_1), (L_2, \phi_2) \in
\mathrm{CY}^{\mathrm{Siegel\text{-}aut}}_2$, there is a natural
Reshetikhin--Turaev comparison map
\[
  \mu_{L_1, L_2}\colon
  \Psi(L_1, \phi_1) \otimes \Psi(L_2, \phi_2)
  \;\xrightarrow{\;\simeq\;}\;
  \Psi\bigl((L_1, \phi_1) \oplus (L_2, \phi_2)\bigr),
\]
together with a unit map $\eta: \mathbb{1}_{\mathrm{QHopf}} \to
\Psi(\mathbf{0}, 1)$, satisfying the lax monoidal coherence diagrams
(associativity pentagon in Lurie HA 5.3.2.1 and the unit triangle) up
to canonical $A_\infty$-gauge.  The $\mu_{L_1,L_2}$ are
quasi-isomorphisms on the Koszul locus of the programme and maps of
quasi-Hopf super-algebras outside it.
\end{theorem}

\begin{proof}[Proof sketch]
Functoriality on morphisms is the Borcherds-pullback compatibility of
Borcherds~1998 \S4 Lemma~4.1 (the singular theta correspondence commutes
with lattice embeddings), transported to the chiral side through
Clause~(iii) of Theorem~\ref{thm:phi-universal-trace} (the
$\Phi$-pushforward of the Koszul reflection).  The comparison map
$\mu_{L_1, L_2}$ is constructed from two independent inputs that agree
on the overlap: \emph{(a)} the Gritsenko--Nikulin~1997--1998
exponential-lift K\"unneth identity
$\phi_1 \boxtimes \phi_2 = \mathrm{Grit}_{L_1 \oplus L_2}(\phi_1 \boxtimes \phi_2)$
when restricted to the sublocus where both lifts exist; \emph{(b)} the
Reshetikhin--Turaev~1991 tensor-product rule for quasi-triangular
quasi-Hopf algebras, stating
$(H_1 \otimes H_2, R_1 \otimes R_2, \Phi_1 \otimes \Phi_2)$ is itself a
quasi-triangular quasi-Hopf algebra with the product $R$-matrix.  The
coherence of $\mu$ under lax symmetry follows from Fubini on the
Jacobi-form indices and the pentagon for $\otimes$ on
$\mathrm{QHopf}^{\mathrm{BKM}}$ (Lurie HA~5.3.2.1 applied to the
symmetric monoidal $\infty$-category of $A_\infty$-quasi-Hopf
super-algebras).  The lax-not-strong qualifier arises because on the
non-Koszul locus the Etingof--Kazhdan gauge class can obstruct strict
equality; the obstruction lives in the $\mathbb{Z}/8$ Bruinier class
of Theorem~\ref{thm:qca-hbar-minus-eighth-lusztig}, vanishing on the
Koszul locus.  For detail on the lax-strong obstruction see Lurie HA
6.3.1.10 (lax monoidal = left-lax from the operadic viewpoint);
Francis~2013 \S5 for the $E_n$-monoidal variant used here.
\end{proof}

\subsection{Test case: \texorpdfstring{$\Psi(\mathrm{II}_{1,1} \oplus \mathrm{II}_{1,1})$}{Psi(II-1,1 (+) II-1,1)}}
\label{subsec:phiutp-test-two-hyperbolic}

\begin{proposition}[Rank-\texorpdfstring{$4$}{4} BKM on \texorpdfstring{$\mathrm{II}_{2,2}$}{II-2,2}]
\label{prop:phiutp-test-rank-four}
\ClaimStatusProvedHere
Let $\mathrm{II}_{1,1}$ be the standard hyperbolic plane lattice and
$\phi_{\mathrm{Heis}}$ the Monster-Heisenberg Jacobi input
\textup{(}Borcherds 1992 Invent.\ Math.\ 109; rank-$1$ Lorentzian,
lift $J(\tau) - 744$\textup{).}  Then
\[
  \Psi\bigl(\mathrm{II}_{1,1} \oplus \mathrm{II}_{1,1},\,
  \phi_{\mathrm{Heis}} \boxtimes \phi_{\mathrm{Heis}}\bigr)
  \;\simeq\;
  \mathbf{H}_{\mathrm{Monster}} \otimes \mathbf{H}_{\mathrm{Monster}},
\]
a rank-$4$ quantum chiral algebra on the even unimodular lattice
$\mathrm{II}_{2,2} = \mathrm{II}_{1,1} \oplus \mathrm{II}_{1,1}$, with
universal $R$-matrix $R_{11} \otimes R_{22}$ and the Drinfeld associator
product $\Phi_1 \otimes \Phi_2$.  The comparison map $\mu$ is a
quasi-isomorphism of quasi-Hopf super-algebras.
\end{proposition}

\begin{proof}[Proof sketch]
The Borcherds lift of $\phi_{\mathrm{Heis}} \boxtimes \phi_{\mathrm{Heis}}$
on $\mathrm{II}_{2,2}$ factorises as a product of two Borcherds lifts on
$\mathrm{II}_{1,1}$ by Borcherds~1998 \S14 K\"unneth
(singular theta correspondences multiply under lattice direct sum).  The
chiral-side factor $\mathbf{H}_{\mathrm{Monster}}$ is the Monster module
vertex algebra $V^\natural$ of Frenkel--Lepowsky--Meurman~1988, realised
on each hyperbolic plane factor as the $E_1$-chiral algebra with
character $J(\tau) - 744$.  The tensor-product rule
$R_{11} \otimes R_{22}$ is the Reshetikhin--Turaev~1991 Proposition~1.5
rule on quasi-triangular quasi-Hopf algebras.  The associator product
is the pentagon for $\otimes$ (Drinfeld~1989 \S3.3, transported to the
BKM setting by Etingof--Kazhdan~2007 Part~V).  Quasi-isomorphism of
$\mu$ follows because both sides have the same
$\mathrm{II}_{2,2}$-lattice Heisenberg free-field content at rank~$1$
Lorentzian $\times$ rank~$1$ Lorentzian, with the only non-trivial
gauge ambiguity in the Drinfeld associator factor, and that ambiguity
is absorbed into the $\mathbb{Z}/8$ Bruinier Chern class which vanishes
on the unimodular $\mathrm{II}_{2,2}$ domain (Bruinier~2002 \S5.1,
applied at discriminant $1$).
\end{proof}

\begin{remark}[From rank-\texorpdfstring{$2$}{2} BKM at \texorpdfstring{$\mathrm{II}_{1,1}$}{II-1,1} to rank-\texorpdfstring{$4$}{4} BKM at \texorpdfstring{$\mathrm{II}_{2,2}$}{II-2,2}]
\label{rem:phiutp-rank-doubling}
The test case witnesses that lattice direct sum on the input side is
not merely set-theoretic doubling: it is a functorial operadic
composition, compatible with the universal $R$-matrix tensor product
and the coproduct factorisation on the output side.  The rank-$4$
BKM algebra $\mathbf{H}_{\mathrm{Monster}} \otimes \mathbf{H}_{\mathrm{Monster}}$
is not a new object -- it is forced by the operadic structure of $\Psi$
together with Borcherds~1998 \S14 K\"unneth.  In particular, the test
case shows that $\Psi$ respects not just composition (functoriality on
morphisms) but also parallel composition (tensor of objects); this is
what elevates $\Psi$ from ``a rule'' to ``a functor of operads''.
\end{remark}

\subsection{\texorpdfstring{$E_n$}{E-n}-structure: \texorpdfstring{$E_2 \to E_1$}{E-2 -> E-1} Dunn stabilisation}
\label{subsec:phiutp-en-structure}

\begin{theorem}[\texorpdfstring{$\Psi$}{Psi} is an \texorpdfstring{$E_2 \to E_1$}{E-2 -> E-1} Dunn stabilisation]
\label{thm:phiutp-e2-to-e1-dunn}
\ClaimStatusProvedElsewhere
As a functor of symmetric monoidal $(\infty,1)$-categories:
\begin{enumerate}[label=\textup{(\roman*)}]
  \item $(\mathrm{CY}^{\mathrm{Siegel\text{-}aut}}_2, \oplus)$
        carries a natural $E_2$-structure: the lattice direct-sum
        pairing $\oplus$ and the Jacobi-form exterior tensor product
        $\boxtimes$ commute by Fubini, giving by Dunn additivity
        $E_1 \otimes_{\mathrm{BV}} E_1 \simeq E_2$
        \textup{(}Dunn~1988 Theorem~3.4; Lurie HA~5.1.2.2\textup{).}
  \item $(\mathrm{QHopf}^{\mathrm{BKM}}, \otimes)$ carries a natural
        $E_1$-structure: the algebra tensor product $\otimes$ is a
        single associative pairing on the $(\infty,1)$-category of
        $A_\infty$-quasi-Hopf super-algebras; iteration does not
        produce a second commuting pairing (the coproduct does not give
        a monoidal structure on $\mathrm{QHopf}$ -- it is an
        \emph{internal} structure on each object).
  \item $\Psi$ is $E_2 \to E_1$ Dunn-stable: the $E_2$-structure on the
        domain projects onto the $E_1$-structure on the codomain by
        forgetting one of the two commuting pairings (the
        Jacobi-form $\boxtimes$ is absorbed into the quasi-Hopf
        coproduct structure internal to the BKM output).
\end{enumerate}
\end{theorem}

\begin{proof}[Proof sketch]
Clause~(i) is Dunn additivity at the monoidal-category level: when a
symmetric monoidal category carries two compatible pairings that commute
up to canonical isomorphism, Dunn's theorem
(Dunn~1988 Theorem~3.4; Lurie HA~5.1.2.2) identifies the iterated
structure with $E_2$.  Clause~(ii) is the standard observation that the
$(\infty,1)$-category of associative algebras (equivalently:
$E_1$-algebras) carries only an $E_1$-structure, not an $E_2$-structure,
because the Eckmann--Hilton argument would then collapse the algebra to
commutative
\textup{(}Lurie HA~5.1.2.9: $E_{n+1}$-algebras in $E_m$-categories are
$E_{n+m}$-algebras; for $H$ itself to support a second pairing compatible
with its multiplication would force $H$ commutative, contradicting
non-abelian BKM\textup{).}  Clause~(iii) is the Dunn-stabilisation
compatibility: the forgetful functor from $E_2$-algebras to
$E_1$-algebras remembers only one pairing, and the preservation of
that forgetful on $\Psi$ is exactly the lax symmetric-monoidal
structure from Theorem~\ref{thm:phiutp-psi-functor-lax-symmon}
(the map $\mu$ is the natural $E_2 \to E_1$ rigidification witness).
\end{proof}

\subsection{Reshetikhin--Turaev compatibility: disjoint union = multiplicative}
\label{subsec:phiutp-rt-compatibility}

\begin{corollary}[Reshetikhin--Turaev multiplicativity for \texorpdfstring{$\Psi$}{Psi}]
\label{cor:phiutp-rt-multiplicative}
\ClaimStatusProvedHere
For a disjoint union of chiral sites $\Sigma = \Sigma_1 \sqcup \Sigma_2$
hosting Borcherds data $(L_1, \phi_1)$ on $\Sigma_1$ and $(L_2, \phi_2)$
on $\Sigma_2$, the quantum invariant
$Z_\Psi(\Sigma) := Z_{\mathrm{RT}}(\Sigma; \Psi(L, \phi))$ with
$(L, \phi) = (L_1, \phi_1) \oplus (L_2, \phi_2)$ factorises
multiplicatively:
\[
  Z_\Psi(\Sigma_1 \sqcup \Sigma_2)
  \;=\;
  Z_\Psi(\Sigma_1) \cdot Z_\Psi(\Sigma_2).
\]
Equivalently, $\Psi$ followed by the Reshetikhin--Turaev functor
$Z_{\mathrm{RT}}: \mathrm{QHopf}^{\mathrm{BKM}} \to \mathrm{TQFT}_3$
\textup{(}Reshetikhin--Turaev 1991, Turaev 1994\textup{)} is a symmetric
monoidal functor from
$(\mathrm{CY}^{\mathrm{Siegel\text{-}aut}}_2, \oplus)$ into the
symmetric monoidal category of $3$-TQFTs with disjoint-union pairing.
\end{corollary}

\begin{proof}
$Z_{\mathrm{RT}}$ is symmetric monoidal
\textup{(}Turaev~1994 Theorem III.4.5: the Reshetikhin--Turaev
invariant of a disjoint union is the product of the invariants\textup{).}
The composite $Z_{\mathrm{RT}} \circ \Psi$ sends
$\oplus$ to $\otimes$ (via $\Psi$) and then $\otimes$ to $\cdot$ (via
$Z_{\mathrm{RT}}$), giving multiplicativity on disjoint unions.  The
coherence (associativity, unit, symmetry) follows from composing the
lax symmetric-monoidal structure on $\Psi$ from
Theorem~\ref{thm:phiutp-psi-functor-lax-symmon} with the strict symmetric
monoidal structure on $Z_{\mathrm{RT}}$ \textup{(}Turaev~1994
Theorem~III.4.6\textup{).}
\end{proof}

\begin{remark}[The four pillars of the \texorpdfstring{$\Psi$}{Psi} operadic upgrade]
\label{rem:phiutp-four-pillars}
The operadic upgrade of $\Psi$ rests on four citations that together
form a proof skeleton: Borcherds~1998 \S14 (K\"unneth for singular
theta correspondence on lattice direct sums), Dunn~1988 Theorem~3.4
(additivity of $E_n$ under Boardman--Vogt tensor product),
Reshetikhin--Turaev~1991 Proposition~1.5 (tensor-product rule for
quasi-triangular quasi-Hopf algebras), and Lurie HA~5.3.2.1 (coherence
of lax monoidal functors in symmetric monoidal
$(\infty,1)$-categories).  The four pillars are independent and each is
a primary citation in its respective area; together they give
$\Psi: (\mathrm{CY}^{\mathrm{Siegel\text{-}aut}}_2, \oplus) \to
(\mathrm{QHopf}^{\mathrm{BKM}}, \otimes)$ as a genuine functor of
operads, not merely as a rule on objects.  Cross-reference: Vol~II
Chapter on factorization algebras (Costello--Gwilliam framework),
where the chiral-side counterpart of $\Psi$ -- the factorization
functor on disjoint chiral sites -- is constructed directly on
$E_2$-chiral algebras.
\end{remark}

\section{The \texorpdfstring{$5$}{5}d holomorphic Chern--Simons realisation of the universal trace}
\label{sec:phiutp-5d-hcs-cfg-realisation}

The universal $\Phi$-trace identity of
Theorem~\ref{thm:phi-universal-trace} admits a direct chain-level BV
realisation as the partition-function pairing of $5$d holomorphic
Chern--Simons theory on $Y = \R_t\times\C^2$ with gauge algebra
$\frakg$. This section records the realisation with explicit BV
field content, propagator, Feynman-graph expansion, one-loop anomaly
split, and all-orders convergence; it traces the universal trace
through the BV-cohomological pairing that Costello--Yagi
\textup{(}arXiv:1810.01970,~\cite{CostelloYagi2018}\textup{)} and
Costello \textup{(}arXiv:1303.2632,~\cite{Costello2013}\textup{)}
establish for simply-laced $\frakg$. The $3$d topological Chern--Simons
companion of Costello--Francis--Gwilliam \textup{(}2026,
\cite{CFG2026}\textup{)} is the Wick-rotated topological boundary of
this construction; the $\R_t$-direction of $5$d hCS is the CFG 2026
topological factor, and the $\C^2$-direction is its holomorphic
enhancement.

\subsection{The BV field content}
\label{subsec:phiutp-5d-hcs-bv-fields}

\begin{definition}[Off-shell BV field content of \texorpdfstring{$5$}{5}d hCS]
\label{def:phiutp-5d-hcs-bv-fields}
Let $\frakg$ be a finite-dimensional complex Lie algebra with
non-degenerate invariant pairing $\mathrm{tr}_\frakg$. The off-shell
BV field content of $5$d holomorphic Chern--Simons on
$Y = \R_t\times\C^2$ with gauge $\frakg$ is the graded vector space
\[
  \mathrm{Fields}_{\mathrm{BV}}(Y, \frakg)
  \;=\; \Omega^{0,\bullet}(\C^2)\otimes \Omega^\bullet(\R_t)\otimes \frakg[1],
\]
with BV differential
$Q_{\mathrm{BV}} = d_t + \bar\partial_{\C^2} + \mathrm{ad}_{\cA}$
and BV antibracket of degree $-1$
\[
  \{-,-\}_{\mathrm{BV}}(\alpha, \beta)
  \;=\; \int_{Y}\, \mathrm{tr}_\frakg(\alpha\wedge\beta)\,
  \wedge\, \Omega_{\C^2},
\]
paired against the flat holomorphic top-form
$\Omega_{\C^2} = dz^1\wedge dz^2$. The classical action is
\begin{equation}
\label{eq:phiutp-5d-hCS-action}
  S_{5\mathrm{d\,hCS}}[\cA]
  \;=\; \int_{Y}\, \mathrm{tr}_\frakg
    \Bigl(\tfrac{1}{2}\cA\wedge(\bar\partial_{\C^2}+d_t)\cA
         + \tfrac{1}{6}\cA\wedge [\cA, \cA]\Bigr)\, \Omega_{\C^2}
\end{equation}
with $\cA$ the superfield running over
$\mathrm{Fields}_{\mathrm{BV}}(Y, \frakg)$. The classical master equation
$\{S, S\}_{\mathrm{BV}} = 0$ holds by the Jacobi identity on $\frakg$
and the Bianchi identity $\bar\partial^2 = 0$.
\end{definition}

\begin{proposition}[Stage-\texorpdfstring{$1$}{1} realisation at \texorpdfstring{$\frakg = \widehat{\mathfrak{gl}}_1$}{g = affine-gl1}]
\label{prop:phiutp-stage1-5d-hcs-gl1}
\ClaimStatusProvedElsewhere
For $\frakg = \widehat{\mathfrak{gl}}_1$, the chain-level Stage-$1$
output $\PhiFA_3(\Perf(\C^3))$ coincides, up to quasi-isomorphism of
$E_3$-factorisation algebras, with the BV observables of $5$d hCS on
$\R_t\times\C^2$ with gauge $\widehat{\mathfrak{gl}}_1$, Wick-rotated
to the topological limit on $\R_t$:
\[
  \PhiFA_3(\Perf(\C^3))
  \;\simeq\;
  \Obs^{\mathrm{BV}}_{5\mathrm{d\,hCS}}\bigl(\R_t\times\C^2;\,
    \widehat{\mathfrak{gl}}_1\bigr)\,
  \otimes_{\widehat{\mathfrak{gl}}_1}\,
  \End(\mathcal{F}_{\mathrm{Fock}}).
\]
\end{proposition}

\begin{proof}[Attribution]
The proof is the combination of three primary inputs. \textup{(i)} The
Costello--Gwilliam holomorphic-topological tensor theorem
\textup{(}\cite{CostelloGwilliam} vol.~2 \S$4.10$\textup{)} gives the
factorisation
$\Obs^{\mathrm{BV}}_{5\mathrm{d\,hCS}}(\R_t\times\C^2,\frakg)
\simeq \Obs^{\mathrm{BV}}_{1\mathrm{d\,top}}(\R_t,\frakg[1])\boxtimes
\Obs^{\mathrm{BV}}_{4\mathrm{d\,hol}}(\C^2,\frakg)$. \textup{(ii)}
The holomorphic factor is the Costello--Li $4$d hCS factorisation
algebra \textup{(}\cite{CostelloLi2020}\textup{),} whose
$E_4$-structure on $\Obs^{\mathrm{BV}}_{4\mathrm{d\,hol}}(\C^2,\frakg)$
is assembled by Kontsevich $E_d$-formality~\cite{Kontsevich1999}
and Tamarkin deformation-quantisation~\cite{Tamarkin2007}. \textup{(iii)}
At $\frakg = \widehat{\mathfrak{gl}}_1$ the Costello--Li $4$d hCS
factorisation algebra is the polyvector-field
Lie-conformal-algebra envelope of $\HH^\bullet(\Perf(\C^3))
\cong \bigoplus_{p+q = \bullet}H^q(\C^3, \wedge^p T_{\C^3})$ which is
precisely the Stage-$1$ output $\PhiFA_3(\Perf(\C^3))$ after
Serre-duality restriction. The compactification on
$\R_t$ corresponds to factorisation homology along the topological
direction and supplies the $E_1$-structure consistent with
Dunn--Lurie additivity $E_4\otimes E_1 \simeq E_5$ and its
$S^2$-framing reduction to $E_3$
\textup{(}Lurie \emph{HA} 5.1.2.2; Costello--Li~\cite{CostelloLi2020}
holomorphic-topological twist\textup{).} See
Proposition~\texttt{prop:5d-hcs-bv-factorisation} of
Chapter~\ref{ch:cy-to-chiral} for the chain-level derivation.
\end{proof}

\subsection{The heat-kernel propagator}
\label{subsec:phiutp-5d-hcs-propagator}

\begin{proposition}[Heat-kernel regularised propagator]
\label{prop:phiutp-5d-hcs-propagator}
The Costello--Li heat-kernel propagator on the $\C^2$-factor is
\[
  P_{\epsilon, L}^{\C^2}
  \;=\; \int_{\epsilon}^{L}\, e^{-t\,\Box_{\C^2}}\,\bar\partial^*\, dt,
\qquad
  \Box_{\C^2}
  \;=\; \bar\partial\bar\partial^* + \bar\partial^*\bar\partial,
\]
with analytic continuation
$P = \lim_{\epsilon\to 0^+}\lim_{L\to\infty}P_{\epsilon, L}
= \bar\partial^*/\Box_{\C^2}$. The full $5$d hCS propagator on
$Y = \R_t\times\C^2$ is the tensor product
\begin{equation}
\label{eq:phiutp-5d-hCS-propagator}
  P_{5\mathrm{d\,hCS}}(t_1, z_1;\, t_2, z_2)
  \;=\; \Theta(t_2 - t_1) \cdot
  \mathrm{BM}(z_1, z_2),
\end{equation}
where $\Theta$ is the Heaviside step kernel on $\R_t$ and
$\mathrm{BM}$ is the Bochner--Martinelli kernel on $\C^2$.
Differences $P_{\epsilon, L} - P_{\epsilon', L'}$ are exact for the
BV differential; the BV cohomology of the observables is independent
of the scheme \textup{(}Costello--Gwilliam~\cite{CostelloGwilliam}
vol.~1 \S$9.3$; Costello--Li~\cite{CostelloLi2016} Proposition~$2.4$\textup{).}
\end{proposition}

\subsection{The Feynman-graph expansion}
\label{subsec:phiutp-5d-hcs-feynman}

\begin{theorem}[Chain-level Feynman expansion of the universal trace]
\label{thm:phiutp-5d-hcs-feynman-universal-trace}
\ClaimStatusProvedHere
The universal trace of Theorem~\ref{thm:phi-universal-trace} admits
the chain-level Feynman-graph expansion
\begin{equation}
\label{eq:phiutp-feynman-expansion}
  \mathrm{tr}_{\mathfrak{Z}(\mathcal{A}_X)}(\mathfrak{K}_{\mathcal{A}_X})
  \;=\;
  \sum_{\Gamma\in \mathrm{Graphs}}\,
  \frac{\hbar^{L(\Gamma)}}{|\mathrm{Aut}(\Gamma)|}\,
  W_\Gamma(P_{5\mathrm{d\,hCS}}, S_{5\mathrm{d\,hCS}})\,
  I_\Gamma(\frakg)\,
  \mathrm{Ch}_\Gamma(\mathfrak{K}),
\end{equation}
where the sum runs over connected trivalent Feynman graphs $\Gamma$
with loop order $L(\Gamma)$; the Kontsevich weight
$W_\Gamma(P, S) = \int_{\mathrm{Conf}(\Gamma, Y)}\prod_{e\in E(\Gamma)}
P_e(x_{s(e)}, x_{t(e)})$ is the configuration-space integral of the
propagator~\eqref{eq:phiutp-5d-hCS-propagator} over the configuration
space of graph-vertices in $Y$; the gauge-invariant
$I_\Gamma(\frakg) = \sum_{\{a_e\}}\prod_{v\in V(\Gamma)} f^{a_1 a_2 a_3}$
contracts the Lie-algebra structure constants along $\Gamma$; and
$\mathrm{Ch}_\Gamma(\mathfrak{K})$ is the Chern character of the
Koszul involution restricted to $\Gamma$.
\end{theorem}

\begin{proof}
The BV partition function $e^{S/\hbar}$ admits the Feynman-graph
expansion by Costello--Gwilliam~\cite{CostelloGwilliam} vol.~$1$ \S$9$
\textup{(}the perturbative BV formalism and its heat-kernel
regularisation\textup{).}
The identification of the partition function with the universal trace
is the chain-level realisation of the Chiral Hochschild Trinity Theorem
\textup{(}Vol~II, \texttt{thm:chiral-hochschild-trinity}\textup{):} the
factorisation-homology output
$\mathfrak{Z}(\mathcal{A}) = \int_{S^1\times\Ran(X)}\mathcal{A}$ is
computed by the BV partition function of $5$d hCS on the corresponding
spacetime, with the $S^1$-direction identified with the
$\R_t$-direction after compactification and the $\Ran(X)$-direction
identified with the chiral observables on the $\C^2$-factor (Lurie
\emph{HA}~$5.5.3.11$: factorisation homology of tensor-product
factorisation algebras is the tensor product of factorisation
homologies). The Koszul involution
$\mathfrak{K}_{\mathcal{A}} : \mathcal{A}\mapsto\mathcal{A}^!$ acts on
the BV field content by $\cA\mapsto -\cA^!$ (Verdier duality pairing);
its Chern character decomposes along the graph expansion into
per-graph contributions. The one-loop correction
$W_{L=1}[P, S]$ reduces to the wheel-diagram sum by
Costello--Yagi~\cite{CostelloYagi2018} Theorem~$4.1$.
\end{proof}

\subsection{The one-loop anomaly: cubic vs.\ quadratic Casimir split}
\label{subsec:phiutp-one-loop-anomaly-split}

\begin{proposition}[Cubic vs.\ quadratic Casimir split of the one-loop anomaly]
\label{prop:phiutp-one-loop-anomaly-split}
\ClaimStatusProvedHere
The one-loop anomaly of the BV partition function
of~\eqref{eq:phiutp-5d-hCS-action} splits into two contributions of
distinct physical character:
\begin{itemize}
  \item \textbf{Scheme-dependent wave-function renormalisation.}
        The quadratic Casimir contribution
        $A_{\mathrm{w.f.}} = -C_2(\frakg)/(2\pi)^3$, with
        $C_2(\frakg) = 2h^\vee$ the dual Coxeter number, arises from
        the propagator self-loop at a single vertex. This term is
        BV-trivial: it is absorbed into the Costello--Li
        heat-kernel regularisation counter-term by adjusting the
        wave-function of the gauge field. It does not obstruct
        quantisation.
  \item \textbf{Cohomological true anomaly.}
        The cubic Casimir contribution
        $A_{\mathrm{anom}} = d^{abc}d^{abc}/(2\pi)^3$, with
        $d^{abc} = \mathrm{str}_\frakg(T^a T^b T^c + T^a T^c T^b)$ the
        totally-symmetric rank-$3$ Casimir, arises from the trivalent
        wheel diagram $W_{\mathrm{wheel}_3}$. This term is
        cohomologically non-trivial in the BV complex in general; its
        vanishing is the gauge-anomaly-freedom condition.
\end{itemize}
For simply-laced $\frakg$ \textup{(}the ADE series\textup{)} the
cyclic-trace identity
$\mathrm{tr}(T^a T^b T^c + T^a T^c T^b) = 0$ forces $d^{abc} = 0$, so
$A_{\mathrm{anom}} = 0$ and the BV quantisation is unobstructed. For
$\mathfrak{sl}_{N\ge 3}$ non-abelian the $d^{abc}$ is non-vanishing;
the theory is obstructed and requires flavour-symmetry cancellation.
For $\widehat{\mathfrak{gl}}_1$ abelian the $d^{abc} = 0$ tautologically;
the theory is unobstructed.
\end{proposition}

\begin{proof}
The one-loop wheel diagram $W_{\mathrm{wheel}_3}$ has three vertices
connected pairwise by three propagators. On the
$\C^2$-factor of $Y$, each propagator contributes the Bochner--Martinelli
kernel; on the $\R_t$-factor, $\Theta$-ordered kernels. The
Kontsevich-weight integral of $W_{\mathrm{wheel}_3}$ over the
configuration space of three vertices on $Y$ evaluates to
$1/(2\pi)^3$ by the standard Bochner--Martinelli integral over
$\mathrm{Conf}_3(\C^2)\times \mathrm{Conf}_3(\R_t) /
\mathrm{Aut}(\mathrm{wheel}_3)$. The Lie-algebra invariant at
$W_{\mathrm{wheel}_3}$ is $d^{abc}d^{abc}$, the totally-symmetric
rank-$3$ Casimir squared. This is the cohomologically true anomaly
of Costello--Yagi~\cite{CostelloYagi2018} Theorem~$4.1$. The
propagator self-loop at a single vertex is a scheme-dependent
logarithmic divergence absorbed by Costello--Li~\cite{CostelloLi2016}
heat-kernel counterterm, contributing $-C_2(\frakg)/(2\pi)^3$ to the
wave-function renormalisation, which is BV-trivial.
\end{proof}

\subsection{All-orders convergence for simply-laced \texorpdfstring{$\frakg$}{g}}
\label{subsec:phiutp-all-orders-convergence}

\begin{theorem}[Costello--Yagi all-orders: simply-laced Yangian VOA]
\label{thm:phiutp-all-orders-costello-yagi}
\ClaimStatusProvedElsewhere
For $\frakg$ simply-laced bosonic \textup{(}$\mathrm{A}_n, \mathrm{D}_n,
\mathrm{E}_6, \mathrm{E}_7, \mathrm{E}_8$\textup{)}, the BV partition
function of $5$d hCS on $\R_t\times\C^2$ quantises to all orders in
$\hbar$ and converges \textup{(}not merely asymptotically\textup{):}
\[
  \bigl\langle e^{S/\hbar}\bigr\rangle_{\mathrm{BV}}
  \;=\; \sum_{n\ge 0}\hbar^n\, Z^{(n)}_{5\mathrm{d\,hCS}}(\frakg),
\qquad
  |Z^{(n)}_{5\mathrm{d\,hCS}}(\frakg)|\; \le\; C^n/n!,
\]
with $C = C(\frakg)$ a positive constant bounded by the
Kontsevich-weight entropy on $E_3$-graph complex. The factorisation
homology along the topological direction $\R_t$ is the Yangian
vertex operator algebra:
\[
  \int_{\R_t}\, \Obs^{\mathrm{BV}}_{5\mathrm{d\,hCS}}(\R_t\times\C^2, \frakg)
  \;\simeq\; Y^{\mathrm{VOA}}(\frakg).
\]
\end{theorem}

\begin{proof}[Attribution]
This is Costello--Yagi~\cite{CostelloYagi2018} Theorem~$4.1$ combined
with the absolute convergence of Kontsevich weights on the
$E_3$-graph complex \textup{(}Kontsevich~\cite{Kontsevich1999}
Theorem~$2$; Brown~\cite{Brown2012}\textup{).} The chain-level
Feynman-graph expansion of
Theorem~\ref{thm:phiutp-5d-hcs-feynman-universal-trace} is the
expansion in question; the convergence bound comes from the
factorial suppression of Kontsevich weights. The identification with
$Y^{\mathrm{VOA}}(\frakg)$ is the content of
Theorem~\texttt{thm:non-abelian-5d-hcs-yangian-voa} of
Chapter~\ref{ch:cy-to-chiral}.
\end{proof}

\subsection{CFG 2026 topological compactification}
\label{subsec:phiutp-cfg-topological-compactification}

\begin{remark}[CFG 2026 as the topological boundary of \texorpdfstring{$5$}{5}d hCS]
\label{rem:phiutp-cfg-2026-topological-boundary}
The Costello--Francis--Gwilliam 2026~\cite{CFG2026} $E_3$-factorisation-algebra
presentation of $3$d ordinary Chern--Simons observables on $\R^3$ is
the topological boundary of the $5$d hCS construction above: compactify
the $\C^2$-factor of $Y = \R_t\times\C^2$ on an $S^2$-fibred
two-sphere and take the Wick-rotated topological limit. On the
Costello--Li~\cite{CostelloLi2020} holomorphic-topological twist side,
this compactification produces the ordinary topological CS observable
algebra on $\R^3$; on the BV partition-function side, it produces the
CFG 2026 factorisation-algebra presentation. The $\R_t$-factor of the
present construction becomes one of the three $\R$-factors of CFG 2026
$\R^3$; the remaining two arise from the topologically-twisted
$\C^2$-factor.

Conversely, the $5$d hCS construction is the holomorphic enhancement
of CFG 2026: along the $\C^2$-direction the observable algebra
carries a holomorphic $E_2$-structure from Kontsevich formality
\cite{Kontsevich1999}, refining the topological $E_3$-structure of
CFG 2026 by an additional holomorphic degree-$0$ layer. Both
algebraic structures survive under $\R_t$-compactification: the
$E_3$-structure becomes the $E_2$-structure on $Y^{\mathrm{VOA}}(\frakg)$
along $\C^2$, and the $E_1$-structure becomes the vertex-algebra
OPE at a point on $\C^2$.

The Universal $\Phi$-Trace Identity, viewed through the $5$d hCS /
CFG 2026 correspondence, is the statement that the BV supertrace of
the Koszul involution $\mathfrak{K}_{\mathcal{A}}$ on the Yangian VOA
$Y^{\mathrm{VOA}}(\frakg)$ equals the Borcherds-lift residue at the
Heegner locus of the Siegel-modular Fourier coefficient
$c_\Lambda(0)/2 = \kappa_{\mathrm{BKM}}(\mathfrak{g}_\Lambda)$. The
two supertraces — Koszul and Borcherds — live on the same $E_3$
factorisation-algebra output at Stage~$1$ of
$\Phi_3 = \SpCh_{\Sigma_2, C}\circ \PhiFA_3$.
\end{remark}

\subsection{Borcherds side from the \texorpdfstring{$5$}{5}d hCS partition function}
\label{subsec:phiutp-borcherds-from-5d-hcs}

\begin{theorem}[Borcherds-lift residue at the Heegner locus of \texorpdfstring{$K3\times E$}{K3 x E}]
\label{thm:phiutp-borcherds-residue-k3e}
\ClaimStatusProvedElsewhere
The Borcherds-side supertrace of Theorem~\ref{thm:phi-universal-trace}
at $X = K3\times E$ is the residue of the $5$d hCS BV partition
function at the rank-$1$ isotropic vector $v\in\Lambda^*$ of the
Mukai lattice $\Lambda = \Lambda_{K3}\oplus \mathrm{II}_{1,1}$:
\[
  \mathrm{tr}_{\mathfrak{Z}(\mathcal{A}_{K3\times E})}(\mathfrak{B}_{K3\times E})
  \;=\; \operatorname*{Res}_{v\cdot v = 0}
  \bigl\langle e^{S/\hbar}\bigr\rangle_{\mathrm{BV},\, \cA|_{C_v}}
  \;=\; c_1(0)/2 \;=\; 5,
\]
with $C_v\subset\C^2$ the curve associated to $v$, $c_1(0) = 10$ the
Fourier coefficient of the weight-$5$ Jacobi form $\phi_{0,1}$ at
$(m, n) = (0, 0)$, and the residue evaluated at the standard
Borcherds singular-theta lift of the BV partition function.
\end{theorem}

\begin{proof}[Attribution]
Clause~(ii) of Theorem~\ref{thm:phi-universal-trace} together with
Borcherds~\cite{Borcherds1998} Theorem~$10.1$ and
Gritsenko--Nikulin~\cite{GritsenkoNikulin1997} Theorem~$2.1$ give the
residue formula. The BV partition function of
Theorem~\ref{thm:phiutp-all-orders-costello-yagi} compactified on
the Narain $\mathrm{II}_{3,19}$-fibre and singular-theta lifted
computes the Gritsenko $\Delta_5$ automorphic form of weight $5$ on
$\Sp_4(\Z)$. The Fourier coefficient $c_1(0) = 10$ and the weight
$5 = c_1(0)/2$ are the content of the universal Borcherds weight
identity $\kappa_{\mathrm{BKM}}(\Phi_N) = c_N(0)/2$ at $N = 1$.
Cross-reference: Chapter~\texttt{cy\_d\_kappa\_stratification},
Theorem~\texttt{thm:borcherds-weight-kappa-BKM-universal}.
\end{proof}

\subsection{CY-A\texorpdfstring{$_3$}{3} from the \texorpdfstring{$5$}{5}d hCS anomaly cancellation at \texorpdfstring{$\widehat{\mathfrak{gl}}_1$}{affine-gl1}}
\label{subsec:phiutp-cya3-anomaly-cancellation}

\begin{corollary}[CY-A\texorpdfstring{$_3$}{3} from anomaly cancellation at abelian gauge]
\label{cor:phiutp-cya3-anomaly-cancellation}
\ClaimStatusProvedElsewhere
The CY-A$_3$ existence-and-$E_1$-rigidity statement of
Theorem~\texttt{thm:cya3-existence-rigidity} of
Chapter~\ref{ch:cy-to-chiral} is equivalent to the vanishing of the
cohomological true anomaly $A_{\mathrm{anom}}$ of
Proposition~\ref{prop:phiutp-one-loop-anomaly-split} for
$\frakg = \widehat{\mathfrak{gl}}_1$, combined with the all-orders
convergence of Theorem~\ref{thm:phiutp-all-orders-costello-yagi}.
The obstruction-tower argument of CY-A$_3$ is, from the $5$d hCS
perspective, the BV cohomology computation that all higher-loop
wheel diagrams cancel recursively by Kontsevich formality on the
$\C^2$-factor; the abelian vanishing $d^{abc} = 0$ at
$\widehat{\mathfrak{gl}}_1$ is the one-loop input.
\end{corollary}

\begin{proof}[Attribution]
At $\frakg = \widehat{\mathfrak{gl}}_1$ the rank-$3$ Casimir is
trivially zero (the Lie algebra is abelian, so all brackets vanish).
The one-loop anomaly $A_{\mathrm{anom}}$ is therefore zero by
Proposition~\ref{prop:phiutp-one-loop-anomaly-split}. The all-orders
cancellation proceeds by Kontsevich formality on the holomorphic
$\C^2$-factor: the $E_2$-graph complex on $\C^2$ is acyclic in
positive loop order for abelian $\frakg$, forcing all higher-loop
obstructions to be cohomologically trivial in BV. The CY-A$_3$
existence-and-$E_1$-rigidity of
Theorem~\texttt{thm:cya3-existence-rigidity} of
Chapter~\ref{ch:cy-to-chiral} is then the chain-level statement that
the BV partition function at $\widehat{\mathfrak{gl}}_1$ is
well-defined to all orders and is $E_1$-rigid in the conjectural
$E_d$-categorical enhancement.
\end{proof}

\subsection{The open K$3$ Yangian at all orders}
\label{subsec:phiutp-super-yangian}

\begin{remark}[K$3$ Yangian at all orders: orthogonal envelope $Y_\hbar(\mathfrak{so}(4, 20))$ and its Hodge-parity $\Z/2$-super-extension]
\label{rem:phiutp-super-yangian-k3}
For a $\Z/2$-graded Lie algebra $\frakg = \frakg_0 \oplus \frakg_1$,
the graded Casimir $d^{abc}$ need not vanish; only the
$\mathfrak{gl}(1|1)$ case is verified to vanish at one loop. The BV
partition function of $5$d hCS with gauge Lie algebra the orthogonal
Mukai-preserving $\mathfrak{so}(4, 20)$ \textup{(}Kac type
$D_{12}$\textup{)}, or its programme-specific Hodge-parity
$\Z/2$-super-extension $\mathfrak{so}(4 \mid 20)$ with $V_+ = \C^4$
and $V_- = \C^{20}$ both carrying symmetric forms restricted from
the Mukai pairing, is the conjectural all-orders K$3$ Yangian
$Y^{\mathrm{VOA}}_\hbar(\mathfrak{so}(4, 20))$ \textup{(}respectively
$Y^{\mathrm{VOA}}_\hbar(\mathfrak{so}(4 \mid 20))$ at the
Hodge-parity-extended level\textup{)}. Kac $\mathfrak{osp}(4 \mid 20)$
is \emph{not} the gauge algebra: its defining form is symplectic on
the $20$-dimensional odd part, incompatible with the symmetric Mukai
pairing. The all-orders quantisation is open: it requires the
Kontsevich formality of Ginzburg--Schedler~\cite{GinzburgSchedler2008}
extended to $E_4$, which is conjectural. The Hodge involution
$\theta: H^*(K3; \Z) \to H^*(K3; \Z)$ decomposes $\widetilde H(K3)$
into $V_+ = H^0 \oplus H^4 \oplus (H^{2,0} \oplus H^{0,2})_\R$ of
dimension $4$ and $V_- = H^{1,1}_{\mathrm{prim}}$ of dimension $20$;
the signature $(4, 20)$ is the Hodge signature of the Mukai pairing.
See Chapter~\ref{ch:k3-yangian} for the
$\mathfrak{so}(4, 20)$-versus-Kac-$\mathfrak{osp}(4 \mid 20)$
trichotomy.
\end{remark}

\subsection{Chain-level completion checkpoint}
\label{subsec:phiutp-chain-level-completion}

Every step of the $5$d hCS realisation above is chain-level explicit:
the BV complex is written in Definition~\ref{def:phiutp-5d-hcs-bv-fields};
the propagator is the heat-kernel regularisation of
Proposition~\ref{prop:phiutp-5d-hcs-propagator}; the Feynman-graph
expansion is
Theorem~\ref{thm:phiutp-5d-hcs-feynman-universal-trace}; the
one-loop anomaly split is
Proposition~\ref{prop:phiutp-one-loop-anomaly-split}; the
simply-laced cancellation and all-orders convergence are
Theorem~\ref{thm:phiutp-all-orders-costello-yagi}; the Borcherds
residue at $K3\times E$ is
Theorem~\ref{thm:phiutp-borcherds-residue-k3e}. The universal
$\Phi$-trace identity of Theorem~\ref{thm:phi-universal-trace} is the
specialisation of this chain-level BV partition function to the
Koszul and Borcherds involutions, viewed through the Stage-$1$
$\PhiFA_3$ realisation of Proposition~\ref{prop:phiutp-stage1-5d-hcs-gl1}
and the Stage-$2$ curve-specialisation $\SpCh_{\Sigma_2, C}$.

\begin{remark}[Positive-geometry grammar at the \texorpdfstring{$5$}{5}d hCS locus]
\label{rem:phiutp-positive-geometry-5d-hcs}
The BV partition function of~\eqref{eq:phiutp-5d-hCS-action} fits the
universal positive-geometry grammar
$Y^+(X) = H^\bullet_{\mathrm{eq}}(\mathcal{M}^+_{\mathrm{eff}}(X),\phi_W)$.
At $X = \C^3$ the effective moduli space is
$\mathcal{M}^+_{\mathrm{eff}}(\C^3) = \mathrm{Hilb}^\bullet(\C^3)$
with superpotential $W = \mathrm{tr}(x[y,z])$; the equivariant
cohomology is the Schiffmann--Vasserot CoHA
\cite{SchiffmannVasserot2013}, and the $5$d hCS BV partition
function on the matrix-model reduction computes precisely this
cohomology with equivariant parameters $(h_1, h_2, h_3)\in\C^3$
subject to the CY$_3$ constraint $h_1 + h_2 + h_3 = 0$.
Schiffmann--Vasserot's CoHA is the positive half
$Y^+(\widehat{\mathfrak{gl}}_1)$ of the affine Yangian obtained as
factorisation homology of
Theorem~\ref{thm:phiutp-all-orders-costello-yagi} along $\R_t$ at
$\frakg = \widehat{\mathfrak{gl}}_1$.
\end{remark}

\subsection{BV-residue base case at \texorpdfstring{$X = \C^3$}{X = C^3}: the empty-Heegner vanishing}
\label{subsec:phiutp-bv-c3-base-case}

The universal identity $\kBKM(\Phi_N) = c_N(0)/2$ attains its
\emph{base-case} value at the non-compact $\C^3$: the Fourier
coefficient is zero, the Heegner locus empty, the one-loop anomaly
absent, the factorisation \v{C}ech descent trivial. All three
conditions vanish simultaneously, giving $\kBKM^{\mathrm{BV}}(\C^3)
= 0$. This is the terminal fibre against which the Gritsenko $\Delta_5$
interior value $\kBKM^{\mathrm{BV}}(K3\times E) = 5$ of
Theorem~\ref{thm:phiutp-borcherds-residue-k3e} is the first
non-trivial excitation.

\begin{theorem}[Empty-Heegner vanishing of the BV residue at \texorpdfstring{$\C^3$}{C^3}]
\label{thm:bl-bv-c3-vanishing}
\ClaimStatusProvedHere
For $X = \C^3$ with gauge algebra $\frakg = \widehat{\mathfrak{gl}}_1$
on $Y = \R_t\times\C^2$, the Path-$2$ BV-residue specialisation of
the universal $\Phi$-trace identity vanishes identically:
\[
  \kBKM^{\mathrm{BV}}(\C^3)
  \;:=\;
  \operatorname*{Res}_{v\cdot v = 0}
  \bigl\langle e^{S_{5\mathrm{d\,hCS}}/\hbar}\bigr\rangle_\BV^{\widehat{\mathfrak{gl}}_1}\Big|_{X = \C^3}
  \;=\; 0.
\]
The vanishing has three independent chain-level causes, each
surviving the other two as perturbations:
\begin{enumerate}[label=\textup{(\alph*)}]
\item \textbf{Cohomological anomaly zero.} The rank-$3$ Casimir
$d^{abc}d^{abc}/(2\pi)^3 = 0$ at the abelian gauge algebra
$\widehat{\mathfrak{gl}}_1$ (since $f^{abc} = 0$); the BV
quantisation is unobstructed to all orders by the abelian
specialisation of
Proposition~\ref{prop:phiutp-one-loop-anomaly-split}.
\item \textbf{Heegner locus empty.} The Mukai lattice of $\C^3$ is
rank-$1$ degenerate, $\Lambda_{\C^3} = \Z\cdot[\mathrm{pt}]$ with
$\langle[\mathrm{pt}],[\mathrm{pt}]\rangle = 0$; the Borcherds
singular-theta lift $\Phi_\Lambda(\phi)$ of
Borcherds~\cite{Borcherds1998} Theorem~$10.1$ requires signature
$(2, n)$ with $n\geq 1$, so the hypothesis fails and the lift is
identically zero.
\item \textbf{Factorisation descent trivial.} On affine $\C^3$ the
Costello--Gwilliam factorisation-algebra \v{C}ech cohomology
$\mathrm{R}\Gamma(\C^3, \Obs^{\mathrm{BV}}_{5\mathrm{d\,hCS}}(-;\widehat{\mathfrak{gl}}_1))$
reduces to the global section at $U_0 = \C^3$; all higher
local-to-global obstructions vanish by affineness.
\end{enumerate}
Combined, these give the empty-Heegner terminal fibre of the
universal $c_N(0)/2$-family at the formal index $N = 0$.
\end{theorem}

\begin{proof}
Condition~(a) is the abelian specialisation of
Proposition~\ref{prop:phiutp-one-loop-anomaly-split}: the
cohomological contribution $A_{\mathrm{anom}} = d^{abc}d^{abc}/(2\pi)^3$
vanishes at $\widehat{\mathfrak{gl}}_1$ tautologically, and the
wave-function contribution $A_{\mathrm{w.f.}} = -C_2(\widehat{\mathfrak{gl}}_1)/(2\pi)^3 = 0$
by the vanishing of the abelian dual Coxeter number. The BV
partition function collapses to the Gaussian determinant
$\det(d_t + \bar\partial_{\C^2})^{-1/2}$ of the quadratic action
(the cubic interaction $\tfrac{1}{6}\mathrm{tr}([\cA,\cA]\cA)$ is
zero for abelian $\frakg$).

Condition~(b): the Mukai lattice $H^{\mathrm{even}}(\C^3;\Z)$ is
rank $1$, generated by the point class $[\mathrm{pt}]$, with
Mukai pairing $\langle[\mathrm{pt}],[\mathrm{pt}]\rangle_\mathrm{Mukai} = 0$
by the standard self-intersection computation for points on a
smooth variety of dimension $\geq 1$. The Borcherds lift
$\Phi_\Lambda(\phi) : M_k^!(\Gamma)\to\mathcal{F}(\Lambda)$
(Borcherds~\cite{Borcherds1998} Thm.~$10.1$) is defined only for
signature $(2, n)$ lattices with $n\geq 1$; at $\Lambda = \Lambda_{\C^3}$
of signature $(0, 1)$ degenerate the hypothesis fails. The residue
over the empty Heegner locus is zero by empty integration.

Condition~(c): the affine open $U = \C^3$ admits the trivial good
cover $\{U\}$ whose \v{C}ech nerve is a point; the local-to-global
descent
$\mathrm{R}\Gamma(\C^3, \Obs^\BV) = \Obs^\BV(\C^3)$
is the global section itself. Higher \v{C}ech cohomology vanishes
by affineness (Hartshorne~\cite{Hartshorne} III.$3$.$5$ for coherent
sheaves; Costello--Gwilliam vol.\ $1$ \S$6$.$4$--$6$ for
factorisation algebras satisfying the Weiss-cover descent). The
Path-$2$ residue is the local observable at a single point; there
is no \v{C}ech-assembly contribution.

The three conditions are independent: (a) is a Lie-algebraic input,
(b) is a lattice-theoretic input, (c) is a topological input. Their
simultaneous vanishing makes the $\C^3$ base case robust under all
three independent perturbations (non-abelian gauge, non-degenerate
lattice, compact topology); each perturbation, turned on one at a
time, produces a distinct non-trivial $\kBKM$ value, and the three
combined produce the full CHL ladder
$(\kBKM)_{N\in\{1,2,3,4,6\}} = (5, 4, 3, 2, 1)$ on the Family~(i)
scope or $(5, 2, 1, 1, 1)$ on the Family~(ii) twined scope of
Remark~\ref{rem:phi-trace-bkm-two-scopes}.
\end{proof}

\begin{remark}[Three-path consistency at the \texorpdfstring{$\C^3$}{C^3} base case]
\label{rem:phiutp-three-path-base-case}
The base-case value $\kBKM^{\mathrm{BV}}(\C^3) = 0$ is independently
produced by all three proof paths of the universal $\Phi$-trace
identity, each by a structurally distinct mechanism:
\begin{itemize}
\item \emph{Path~$1$ (Hodge supertrace).} The Hodge-theoretic
$\kch(X) = \sum_q(-1)^q h^{0, q}(X)$ evaluated on $X = \C^3$ via
the affine vanishing theorem $h^{0, q}(\C^3) = 0$ for $q > 0$
(Hartshorne~\cite{Hartshorne} III.$3$.$7$). The reduced-cohomology
supertrace matches the Mukai empty-Heegner convention.
\item \emph{Path~$2$ (BV residue).} This theorem: the Borcherds
singular-theta integral over the empty Heegner locus of the
degenerate rank-$1$ Mukai lattice is zero, with additional
supporting vanishings from the abelian anomaly freedom and the
affine factorisation descent.
\item \emph{Path~$3$ (Maulik--Okounkov stable envelope).} The MO
$R$-matrix on the free theory at abelian gauge is the identity
(no wall-crossing residues at $\widehat{\mathfrak{gl}}_1$); the
gluing-cocycle residue
$R^{MO}(u) = \mathrm{Res}_{u = u_\star}\phi^+_{\mathrm{UV}}(u)$
collapses to zero because $\phi^+_{\mathrm{UV}}$ has no poles in
the abelian chamber stratification.
\end{itemize}
The three-path consistency at $\C^3$ is tautological because each
path's vanishing follows from a distinct structural boundary
(affine vanishing in Hodge, empty Heegner in BV, trivial
$R$-matrix in MO). Non-trivial three-path consistency is forced at
the $K3\times E$ interior (Theorem~\ref{thm:phiutp-borcherds-residue-k3e}),
where all three paths deliver the shared value $5$.
\end{remark}

\begin{remark}[Schiffmann--Vasserot CoHA side of the base case]
\label{rem:phiutp-sv-coha-base-case}
The Schiffmann--Vasserot identification
$\CoHA(\C^3) = Y^+(\widehat{\mathfrak{gl}}_1)$ is the CoHA-side
companion to the BV-residue base case. The Yangian
$Y^+(\widehat{\mathfrak{gl}}_1)$ is graded by mode level $n\geq 0$;
the vacuum sector $n = 0$ has Borcherds weight zero (matching
$\kBKM^{\mathrm{BV}}(\C^3) = 0$), and the first-mode generator
$e_1(z)$ has non-trivial Borcherds weight which, after
specialisation to the Gritsenko $\Delta_5$ Heegner lift on
$K3\times E$, is the interior value $\kBKM^{\mathrm{BV}}(K3\times E)
= 5$. The $\C^3$ base case is therefore the vacuum contribution
to the full Yangian-Borcherds-weight spectrum; the $K3\times E$
interior values are the first-mode-and-higher excitations. The
positive-geometry grammar
$Y^+(X) = H^\bullet_{\mathrm{eq}}(\mathcal{M}^+_{\mathrm{eff}}(X), \phi_W)$
of Remark~\ref{rem:phiutp-positive-geometry-5d-hcs} at $X = \C^3$
reads $\mathcal{M}^+_{\mathrm{eff}}(\C^3) = \Hilb^\bullet(\C^3)$,
$W = \mathrm{tr}(x[y, z])$; the base-case Borcherds weight is the
$n = 0$ piece of this graded object.
\end{remark}

\begin{remark}[How non-zero $\kBKM$ emerges from the base case]
\label{rem:phiutp-bv-c3-deformation}
The route from $\kBKM^{\mathrm{BV}}(\C^3) = 0$ to
$\kBKM^{\mathrm{BV}}(K3\times E) = 5$ traverses two independent
deformations.
(i)~\emph{Compactification of the holomorphic base.} The base $\C^2$
is compactified to $K3$ with signature-$(3, 19)$ Mukai lattice; the
period domain acquires signature $(2, 19)$ Heegner divisors, and
the Borcherds lift produces the genuine Gritsenko $\Delta_5$
automorphic form of weight $5$. The Bochner--Martinelli kernel
is replaced by the hyperkähler heat kernel on $K3$; wheel integrals
pick up residues at the $(-2)$-curves of $K3$.
(ii)~\emph{Elliptic direction.} The topological $\R_t$ is replaced
by an elliptic curve $E$; the partition function becomes a trace
on a Hilbert space indexed by the $\Z$-periodicity of $E$,
producing the modular index which exposes the CHL slice
$N\in\{1, 2, 3, 4, 6\}$ via the twined Jacobi forms
$\phi_{0, 1}^{(g_N)}$ of
Gritsenko--Nikulin~\cite{GritsenkoNikulin1997} Theorem~$2.1$.
The simultaneous turning-on of both deformations produces the
five-value CHL ladder $(5, 2, 1, 1, 1)$ on the Family-(ii) twined
scope of Remark~\ref{rem:phi-trace-bkm-two-scopes}; the base case
$\C^3\mapsto 0$ is the zero-extension of this ladder to the formal
index $N = 0$.
\end{remark}
